% BAB III METODE PENELITIAN
\chapter{METODE PENELITIAN}
\label{chap:metode}

Bab ini menjelaskan metode penelitian, yang meliputi pendekatan dan tahapan penelitian, kajian awal lingkungan penelitian, data penelitian, dekomposisi persoalan, kebutuhan awal artefak, pemilihan pendekatan, kerangka berpikir, rancangan artefak, rancangan eksperimen dan evaluasi, pemetaan pertanyaan penelitian, pembagian peran, etika penelitian, alat dan bahan, pengelolaan risiko, rencana luaran dan publikasi, serta jadwal penelitian.

Proposal ini diajukan sebelum tahap Perencanaan Penelitian. Karena itu, Subbab~\ref{sec:lingkungan} sampai dengan Subbab~\ref{sec:kerangka} memuat kajian awal yang bersumber dari studi dokumen dan wawancara pada 23 Oktober 2026, bukan analisis yang telah lengkap. Profil agregat awal berkas dinas disajikan pada Subbab~\ref{sec:data}. Analisis yang lebih mendalam, yaitu \textit{data profiling} lengkap termasuk pengukuran kualitas data kondisi sebelum, validasi kebutuhan bersama dinas, dan pemodelan proses \textit{as-is} dan \textit{to-be}, dilaksanakan setelah Seminar Proposal dan dilaporkan pada Laporan Tugas Akhir.

\section{Pendekatan Penelitian}
\label{sec:metode}

Penelitian ini menggunakan paradigma \textit{design science research} (DSR) dengan mengikuti DSRM Peffers dkk. \cite{peffers2007dsrm} (Gambar~\ref{fig:dsr}) dan pedoman Hevner dkk. \cite{hevner2004design}. Kontribusi utama penelitian ini adalah evaluasi empiris efektivitas intervensi, yaitu \textit{prototype} yang dirancang berdasarkan teori dan temuan lapangan; klaim pengetahuan tidak diajukan pada tingkat \textit{prototype} \cite{gregor2013positioning}. Dasar teori rancangan artefak meliputi literatur LLM dan \textit{retrieval-augmented generation} \cite{gao2023ragsurvey,ji2023hallucination}, pembagian peran antara manusia dan mesin \cite{raisch2021automation} serta temuan bahwa meminta pengguna menilai saran AI secara aktif dapat mengurangi ketergantungan pada saran yang keliru \cite{bucinca2021trust}, serta literatur kualitas data \cite{iso25012,govuk2021dqdimensions}, sedangkan dasar lapangannya adalah temuan wawancara pada Tabel~\ref{tab:temuan}. Relevansi penelitian diupayakan melalui keterlibatan dinas dan desa sejak identifikasi masalah hingga evaluasi, sedangkan ketelitian (\textit{rigor}) diupayakan melalui penggunaan landasan teori yang mapan dan evaluasi campuran (\textit{mixed methods}). Evaluasi terdiri atas tiga komponen yang saling melengkapi: (1) eksperimen terkontrol untuk menguji komponen standardisasi KBLI 2025 (RQ1), disertai analisis ablasi eksploratif mengenai konteks bahasa lokal dan analisis sekunder atas pembagian peran manusia dan mesin berbasis skor konsistensi (\textit{coverage--accuracy curve}; Subbab~\ref{sec:eksperimen}), (2) studi kuantitatif sebelum--sesudah untuk mengukur kualitas data pada enam dimensi (RQ2), dan (3) evaluasi lapangan \textit{within-subject} dengan pengukuran kuantitatif dan wawancara pascatugas untuk menilai kinerja tugas dan beban kerja pengguna (RQ3). Pemisahan data pribadi bukan komponen evaluasi, melainkan kebutuhan nonfungsional yang dipenuhi sebagai bagian dari etika penelitian (Subbab~\ref{sec:etika}). Pemeringkatan sasaran pembinaan diimplementasikan sebagai fitur pendukung keputusan yang tidak diuji (Subbab~\ref{sec:ruang-lingkup}). \textit{Prototype} berfungsi sebagai intervensi sekaligus instrumen penelitian; \textit{functional testing} hanya dilakukan untuk memastikan bahwa instrumen tersebut layak digunakan dalam pengumpulan data, bukan sebagai kriteria keberhasilan penelitian.

\textbf{Pemilihan varian DSR.} DSRM Peffers dkk. \cite{peffers2007dsrm} dipakai dengan titik masuk berupa persoalan (\textit{problem-centered initiation}), karena penelitian ini bermula dari persoalan kualitas data yang dilaporkan dinas, bukan dari ketersediaan teknologi tertentu. Dalam kegiatan penelitian, dinas dan desa memvalidasi kebutuhan dan mengikuti uji coba. Tahap evaluasi DSRM menghubungkan eksperimen pengodean pada data uji yang dikunci (RQ1) dengan evaluasi kualitas data dan kinerja operator di desa percontohan (RQ2 dan RQ3). Studi kasus tidak dijadikan pendekatan utama karena fokus penelitian adalah pengembangan dan evaluasi artefak melalui tahapan DSRM yang eksplisit; konteks satu kabupaten tetap digunakan untuk evaluasi lapangan. Eksperimen murni juga tidak dijadikan pendekatan utama karena penelitian mencakup evaluasi pada lingkungan nyata untuk RQ2 dan RQ3.

\textbf{Strategi evaluasi.} Strategi evaluasi dipetakan pada kerangka FEDS \cite{venable2016feds}. Strategi yang dipilih adalah \textit{technical risk and efficacy}, karena risiko utama penelitian ini bersifat teknis, yaitu apakah komponen standardisasi mencapai akurasi yang memadai pada data berbahasa Indonesia, sehingga evaluasi buatan yang ketat didahulukan sebelum evaluasi pada lingkungan nyata yang berbiaya tinggi. Episode evaluasi terdiri atas tiga bagian. Pertama, evaluasi \textit{ex ante} dan formatif berupa lokakarya validasi rancangan dan kebutuhan bersama staf dinas sebelum implementasi (tahap kedua Subbab~\ref{sec:tahapan}). Kedua, evaluasi \textit{ex post} dan sumatif bersifat buatan (\textit{artificial}) berupa eksperimen terkontrol pada \textit{gold standard} untuk RQ1 (Subbab~\ref{sec:eksperimen}). Ketiga, evaluasi \textit{ex post} dan sumatif dalam konteks naturalistik (\textit{naturalistic}) melalui uji coba lapangan di desa percontohan untuk RQ2 dan RQ3 (Subbab~\ref{sec:evalkualitas} dan Subbab~\ref{sec:evalpengguna}); data RQ2 dikumpulkan di lapangan tetapi kualitasnya diukur dan dibandingkan dengan prosedur terkontrol, sedangkan RQ3 melibatkan operator desa dalam tugas terkontrol.

\section{Tahapan Penelitian}
\label{sec:tahapan}

Tahapan penelitian mengikuti enam aktivitas DSRM yang dipetakan ke kegiatan penelitian sebagaimana Gambar~\ref{fig:tahapan}. Aktivitas perancangan dan evaluasi dilaksanakan secara iteratif; apabila hasil evaluasi tidak memenuhi kriteria penerimaan, tim kembali ke aktivitas perancangan.

\begin{figure}[H]
  \centering
  \begin{tikzpicture}[node distance=0.45cm]
  \node (s1) [boxw] {\textbf{1. Identifikasi masalah dan motivasi}\\Studi dokumen LKjIP dan Renstra; wawancara Kadis dan staf (23 Oktober 2026); studi literatur};
  \node (s2) [boxw, below=of s1] {\textbf{2. Penetapan tujuan solusi}\\Kebutuhan fungsional dan nonfungsional; kriteria penerimaan; izin dan akses data};
  \node (s3) [boxw, below=of s2] {\textbf{3. Perancangan dan pengembangan}\\3a. \textit{Data profiling}; \textit{gold standard} KBLI 2025\\3b. \textbf{\textit{Benchmarking} LLM (akurasi, biaya, latensi) dan pemilihan model}\\3c. \textit{Pipeline} standardisasi; modul pendataan; \textit{dashboard} rekapitulasi dan pemeringkatan sederhana};
  \node (s4) [boxw, below=of s3] {\textbf{4. Demonstrasi}\\Uji coba lapangan tiga minggu pada tiga sampai empat desa percontohan bersama staf dinas dan operator desa};
  \node (s5) [boxw, below=of s4] {\textbf{5. Evaluasi}\\Eksperimen klasifikasi (RQ1); kualitas data pada enam dimensi (RQ2); kinerja tugas dan beban kerja pengguna (RQ3)};
  \node (d) [decision, below=0.6cm of s5, text width=2.2cm, aspect=2.2, inner sep=1pt, font=\footnotesize, align=center] {Kriteria\\terpenuhi?};
  \node (s6) [boxw, below=0.6cm of d] {\textbf{6. Komunikasi}\\Laporan Tugas Akhir; Seminar Hasil; Sidang; manuskrip JCIS};
  \draw[arrow] (s1) -- (s2); \draw[arrow] (s2) -- (s3); \draw[arrow] (s3) -- (s4); \draw[arrow] (s4) -- (s5); \draw[arrow] (s5) -- (d);
  \draw[arrow] (d) -- node[right, font=\footnotesize] {Ya} (s6);
  \draw[arrow] (d.east) -- node[below, font=\footnotesize] {Tidak} ++(5.4,0) |- (s3.east);
  \end{tikzpicture}
  \caption{Tahapan Penelitian Berdasarkan DSRM}
  \label{fig:tahapan}
\end{figure}

\begin{enumerate}
  \item \textbf{Identifikasi masalah dan motivasi.} Masukan berupa LKjIP 2025, Renstra 2025--2029, transkrip wawancara, dan literatur. Proses berupa analisis dokumen dan analisis tematik wawancara. Luaran berupa rumusan masalah dan celah penelitian pada Bab~\ref{chap:pendahuluan} serta Subbab~\ref{sec:lingkungan} sampai dengan Subbab~\ref{sec:kerangka}. Tahap ini telah dilaksanakan hingga 23 Oktober 2026.
  \item \textbf{Penetapan tujuan solusi.} Masukan berupa luaran tahap pertama. Proses berupa lokakarya validasi kebutuhan bersama staf dinas dan penetapan kriteria penerimaan. Luaran berupa spesifikasi kebutuhan yang tervalidasi, model proses \textit{as-is} dan \textit{to-be} dalam notasi BPMN, serta surat izin penelitian dan persetujuan akses data.
  \item \textbf{Perancangan dan pengembangan.} Masukan berupa spesifikasi kebutuhan dan data dinas. Proses berupa \textit{data profiling} dan penyusunan \textit{gold standard} KBLI 2025 (3a), \textit{benchmarking} LLM pada data pengembangan untuk memilih model dengan mempertimbangkan akurasi, biaya, dan latensi (3b; Subbab~\ref{sec:benchmark}), serta perancangan basis data dan antarmuka, implementasi \textit{pipeline} standardisasi, serta implementasi \textit{dashboard} rekapitulasi dan pemeringkatan sederhana. Luaran berupa \textit{prototype} sistem yang telah lulus \textit{functional testing}.
  \item \textbf{Demonstrasi.} Masukan berupa \textit{prototype}. Proses berupa pelatihan singkat dan uji coba lapangan selama tiga minggu pada tiga sampai empat desa percontohan (Subbab~\ref{sec:evalpengguna}). Luaran berupa data hasil pendataan melalui sistem dan catatan penggunaan.
  \item \textbf{Evaluasi.} Masukan berupa \textit{gold standard}, data uji coba, dan partisipan. Proses berupa eksperimen dan evaluasi sebagaimana Subbab~\ref{sec:eksperimen} sampai dengan Subbab~\ref{sec:evalpengguna}. Luaran berupa hasil pengukuran untuk RQ1--RQ3.
  \item \textbf{Komunikasi.} Masukan berupa seluruh hasil penelitian. Proses berupa penulisan laporan dan manuskrip. Luaran berupa laporan Tugas Akhir dan manuskrip yang siap diajukan ke JCIS.
\end{enumerate}

\subsection{Fase Benchmarking LLM}
\label{sec:benchmark}

Fase \textit{benchmarking} LLM merupakan bagian aktivitas ketiga DSRM (butir 3b) dan bertujuan memilih LLM untuk \textit{pipeline} standardisasi dengan mempertimbangkan akurasi, biaya, latensi, dan kendala operasional, bukan akurasi saja. Fase ini dilaksanakan pada semester ganjil, setelah \textit{gold standard} pengembangan tersedia dan sebelum titik keputusan G2 (Tabel~\ref{tab:gate}), serta hanya menggunakan 100 \textit{record} pengembangan; 400 \textit{record} uji tidak digunakan. Prosedurnya sebagai berikut.

\begin{enumerate}
  \item \textbf{Penetapan kandidat.} Tiga sampai lima model dipilih, sekurang-kurangnya satu model \textit{open-weight} yang dapat dijalankan pada perangkat keras yang tersedia (Tabel~\ref{tab:tools}) dan satu model \textit{proprietary}. Daftar kandidat, versi, dan tanggal akses dicatat pada \textit{log} eksperimen sebelum pengukuran. Model belum disebutkan pada proposal ini karena ketersediaan dan harga model berubah cepat.
  \item \textbf{Penetapan ambang kelayakan.} Anggaran biaya per seribu \textit{record}, latensi maksimum per \textit{record}, dan kebutuhan menjalankan model secara lokal ditetapkan bersama dinas dan dicatat sebelum pengukuran.
  \item \textbf{Pengukuran.} Seluruh kandidat dijalankan pada kondisi P1 dengan konfigurasi awal yang identik (nilai $k$, \textit{prompt}, dan model \textit{embedding}) pada 100 \textit{record} pengembangan, dengan temperatur nol, menggunakan kriteria pada Tabel~\ref{tab:benchmark}.
  \item \textbf{Aturan keputusan.} Kandidat yang tidak memenuhi syarat kelayakan disisihkan. Pemilihan dilakukan terpisah untuk kelompok \textit{open-weight} dan \textit{proprietary}. Dalam setiap kelompok, kandidat dengan Acc@1 tertinggi menjadi acuan. Jika interval kepercayaan 95\% selisih Acc@1 berpasangan (\textit{bootstrap}) antara kandidat lain dan acuan memuat nol, keduanya masuk tahap pemilihan berdasarkan biaya per seribu \textit{record} terendah, lalu latensi median terendah bila biaya sama. Jika hanya satu kandidat dalam kelompok yang lolos, kandidat itu dipilih. Jika tidak ada yang lolos, tindak lanjut mengikuti keputusan G2 pada Tabel~\ref{tab:gate}. Model \textit{open-weight} terpilih menjadi model utama dan model \textit{proprietary} terpilih menjadi model replikasi (Subbab~\ref{sec:eksperimen}).
\end{enumerate}

Karena hanya 100 \textit{record} yang dipakai, interval kepercayaan antarmodel lebar. Fase ini karenanya berfungsi menyaring kandidat yang layak, bukan membuktikan keunggulan suatu model; perbandingan antarmodel pada data uji tetap dilaporkan sebagai analisis eksploratif. Model tidak diganti setelah data uji dijalankan, dan seluruh data pengukuran serta keputusan dicatat pada \textit{log} eksperimen dan dilaporkan pada Laporan Tugas Akhir.

\section{Kajian Awal Lingkungan Penelitian}
\label{sec:lingkungan}

\subsection{Pemangku Kepentingan dan Proses Bisnis Saat Ini}

Pemangku kepentingan utama penelitian ini adalah Kepala DKUKMPP sebagai pengambil keputusan kebijakan pembinaan, staf teknis Bidang Koperasi dan UMKM sebagai pengelola data dan pelaksana seleksi pembinaan, camat dan perangkat kecamatan sebagai perantara pendataan, serta kepala desa atau operator desa sebagai pihak yang paling dekat dengan pelaku usaha. Pelaku UMKM merupakan subjek data. Proses pendataan saat ini dan proses yang diusulkan digambarkan pada Gambar~\ref{fig:asis-tobe}.

\begin{figure}[H]
  \centering
  \begin{tikzpicture}[node distance=0.55cm and 0.6cm]
  % As-is
  \node[font=\small\bfseries] (t1) {(a) Proses saat ini};
  \node (a1) [box, below=0.3cm of t1] {Dinas menyusun surat permintaan data (ditandatangani Sekda)};
  \node (a2) [box, below=of a1] {Kecamatan meneruskan surat kepada kepala desa};
  \node (a3) [box, below=of a2] {Desa mengisi format lembar kerja atau mengirim data mentah};
  \node (a4) [box, below=of a3] {Kecamatan mengompilasi lalu mengirim berkas ke dinas};
  \node (a5) [boxo, below=of a4] {Dinas menyalin, merapikan, dan merekap secara manual};
  \draw[arrow] (a1) -- (a2); \draw[arrow] (a2) -- (a3); \draw[arrow] (a3) -- (a4); \draw[arrow] (a4) -- (a5);
  % To-be
  \node[font=\small\bfseries, right=3.2cm of t1] (t2) {(b) Proses yang diusulkan};
  \node (b1) [box, below=0.3cm of t2] {Operator desa mengisi atau memutakhirkan data melalui formulir terpandu};
  \node (b2) [box, below=of b1] {Validasi otomatis (format, wilayah, duplikasi)};
  \node (b3) [box, below=of b2] {Standardisasi KBLI dengan \textit{retrieval-grounded} LLM};
  \node (b4) [box, below=of b3] {Verifikasi staf dinas untuk keluaran berskor konsistensi rendah};
  \node (b5) [boxo, below=of b4] {\textit{Dashboard} rekapitulasi dan pemeringkatan sederhana sasaran pembinaan};
  \draw[arrow] (b1) -- (b2); \draw[arrow] (b2) -- (b3); \draw[arrow] (b3) -- (b4); \draw[arrow] (b4) -- (b5);
  \end{tikzpicture}
  \caption{Perbandingan Proses Pendataan UMKM Saat Ini dan Proses yang Diusulkan}
  \label{fig:asis-tobe}
\end{figure}

Proses saat ini memiliki empat titik kelemahan. Pertama, rantai penyampaian data melewati tiga tingkat pemerintahan sehingga tingkat respons rendah dan tidak dapat dipantau; dinas tidak memiliki daftar desa yang telah dan belum mengirim data. Kedua, format isian tidak divalidasi pada saat pengisian sehingga kesalahan baru diketahui ketika data dikompilasi. Ketiga, atribut jenis usaha diisi sebagai teks bebas sehingga rekapitulasi per jenis usaha memerlukan pengodean manual. Keempat, tidak tersedia mekanisme untuk melaporkan usaha yang berhenti beroperasi, sehingga data terus menumpuk tanpa pemutakhiran.

\subsection{Temuan Wawancara dan Implikasi Desain}

Temuan wawancara dengan Kepala Dinas dan staf Bidang Koperasi dan UMKM (\lampkadis{} dan \lampstaf) dipetakan ke implikasi desain artefak pada Tabel~\ref{tab:temuan}.

{\footnotesize
\setlength{\tabcolsep}{4pt}
\renewcommand{\arraystretch}{1.15}
\begin{longtable}{|>{\centering\arraybackslash}p{0.04\textwidth}|L{0.40\textwidth}|L{0.48\textwidth}|}
\caption{Temuan Wawancara Pemangku Kepentingan dan Implikasi Desain}
\label{tab:temuan}\\
\hline
\textbf{No.} & \textbf{Temuan} & \textbf{Implikasi Desain} \\ \hline
\endfirsthead
\hline
\textbf{No.} & \textbf{Temuan} & \textbf{Implikasi Desain} \\ \hline
\endhead
1 & Pemutakhiran data setiap tiga bulan sebagaimana seharusnya tidak pernah terlaksana; data UMKM lama dinas berasal dari pendataan tahun 2020 (sekitar 12.000 unit menurut keterangan lisan) dan tidak digunakan sebagai data penelitian & Modul pemutakhiran berkala dengan status aktif, berhenti, atau pindah, beserta pencatatan tanggal verifikasi untuk mengukur kemutakhiran \\ \hline
2 & Pendataan berjenjang melalui surat Sekda, kecamatan, dan desa dalam bentuk lembar kerja; tingkat respons rendah & Formulir daring per desa dengan akun operator desa dan pemantauan status pengiriman per desa \\ \hline
3 & Data per kecamatan tercampur dan rekapitulasi disusun manual; dinas sering diminta data per kecamatan, klaster, dan jenis usaha & Normalisasi kode wilayah dan KBLI serta rekapitulasi otomatis yang dapat diekspor \\ \hline
4 & Hanya lima personel teknis yang juga menangani Koperasi Desa Merah Putih & Otomatisasi standardisasi dengan antrean verifikasi hanya untuk keluaran yang ditandai perlu verifikasi atau berskor konsistensi rendah \\ \hline
5 & Bantuan dan pelatihan terbatas (misalnya 20 mesin jahit, 30 peserta pelatihan) sehingga dilakukan kurasi dengan kriteria; pemerataan antarwilayah menjadi pertimbangan & Pemeringkatan prioritas sederhana (\textit{Simple Additive Weighting}, SAW) dengan bobot kriteria yang ditetapkan dinas dan batasan kuota per kecamatan; fitur ini tidak diuji \\ \hline
6 & Dinas ingin desa dapat melihat dan memutakhirkan data, tetapi data tidak dibuka untuk umum & Kontrol akses berbasis peran: operator desa hanya mengakses data desanya \\ \hline
7 & Lokasi sebagian desa memerlukan perjalanan sekitar empat jam dari ibu kota kabupaten; Renstra mencatat konektivitas internet yang tidak merata \cite[hlm.~44]{renstra2025toba}
& Antarmuka ringan yang dapat digunakan pada jaringan berkecepatan rendah dengan penyimpanan isian sementara \\ \hline
8 & Pelaku UMKM lebih membutuhkan pendampingan langsung daripada alat digital mandiri; platform pemasaran daerah yang pernah dibangun tidak digunakan & Pengguna sistem adalah aparatur dinas dan desa, bukan pelaku UMKM; sistem berfungsi sebagai augmentasi pendampingan \\ \hline
9 & Instansi provinsi meminta data UMKM untuk dimasukkan ke aplikasi provinsi & Ekspor data dengan struktur atribut baku dan kode referensi \\ \hline
10 & Prioritas pembinaan diberikan kepada wirausaha pemula; sektor yang berkembang adalah tenun ulos dan kriya & Atribut lama usaha dan klaster sektor pembinaan dalam formulir serta kriteria pemeringkatan \\ \hline
\end{longtable}
}

\section{Data Penelitian}
\label{sec:data}

Data yang digunakan mengikuti enam kelompok pada Subbab~\ref{sec:ruang-lingkup} butir 4, yaitu data kiriman desa 2026, data uji coba lapangan, \textit{gold standard}, data referensi, data evaluasi pengguna, serta dokumen dan wawancara pendukung. Data hasil pendataan tahun 2020 tidak digunakan sebagai data penelitian; keterangan mengenai data tersebut hanya menjadi latar masalah (\lampstaf). Subbab ini merinci atribut, tipe data, dan kegunaan kelompok data yang berbentuk tabel. Seluruh contoh nilai pada tabel bersifat fiktif atau tersamarkan dan tidak berasal dari data pelaku usaha.

\subsection{Data Kiriman Desa 2026 (Lembar Kerja Dinas)}
\label{sec:data-dinas}

Data kiriman desa 2026 berupa berkas lembar kerja Microsoft Excel yang memuat daftar UMKM beserta identitas pemilik usaha (nama pengusaha, NIK, alamat, dan nomor telepon), wilayah (desa dan kecamatan), jenis usaha, NIB, jumlah modal usaha, dan keterangan. Berkas yang tersedia bagi tim memuat 10.053 \textit{record}, masing-masing mewakili satu pelaku usaha, dengan nomor urut dan sepuluh atribut sehingga seluruhnya sebelas kolom (Tabel~\ref{tab:dataset}). Lembar kerja tidak memuat omzet maupun tanggal verifikasi; jumlah modal dicatat dalam rupiah, dan kisaran hasil penjualan tahunan baru dihimpun melalui formulir uji coba. Menurut staf, saat wawancara pada Oktober 2026 kiriman desa periode Agustus--September belum selesai dikompilasi (\lampstaf). Akses untuk penggunaan data penelitian diajukan melalui surat resmi kepada dinas, sebagaimana prosedur yang disepakati dalam wawancara (\lampstaf); asal berkas dan tanggal pembaruannya dikonfirmasi kepada dinas pada tahap \textit{data profiling}. Berkas dipakai untuk tiga keperluan, yaitu kerangka sampel \textit{gold standard} (Subbab~\ref{sec:eksperimen}), sumber \textit{record} kondisi sebelum untuk RQ2 (Subbab~\ref{sec:evalkualitas}), dan bahan profil data. Berkas tidak dipakai untuk melatih model apa pun.

Tabel~\ref{tab:dataset} menggunakan nama kolom pada lembar kerja dinas. Tipe data mengikuti konvensi pemrograman umum (\textit{integer}, \textit{float}, \textit{char} dengan panjang tetap, \textit{string}, \textit{boolean}, dan \textit{date}) dan menunjukkan representasi yang direncanakan setelah normalisasi; contoh menggambarkan format nilai. Kolom keterangan memuat hasil profil awal yang dihitung pada 9 November 2026 secara agregat tanpa menampilkan nilai individual. Angka tersebut dapat berubah setelah pembersihan dan dilaporkan lengkap pada Laporan Tugas Akhir. Pada format dinas tidak terdapat atribut bertipe \textit{float}, \textit{boolean}, maupun \textit{date}; atribut bertipe tersebut muncul pada formulir pendataan baru (Tabel~\ref{tab:formulir}) dan pada atribut hasil sistem (Tabel~\ref{tab:atribut-sistem}).

{\footnotesize
\setlength{\tabcolsep}{4pt}
\renewcommand{\arraystretch}{1.15}
\begin{longtable}{|>{\centering\arraybackslash}p{0.04\textwidth}|L{0.13\textwidth}|L{0.10\textwidth}|L{0.21\textwidth}|L{0.40\textwidth}|}
\caption{Atribut Data Kiriman Desa 2026: Tipe Data, Contoh, dan Hasil Profil Awal}
\label{tab:dataset}\\
\hline
\textbf{No.} & \textbf{Atribut} & \textbf{Tipe Data} & \textbf{Contoh (fiktif)} & \textbf{Keterangan dan Hasil Profil Awal} \\ \hline
\endfirsthead
\hline
\textbf{No.} & \textbf{Atribut} & \textbf{Tipe Data} & \textbf{Contoh (fiktif)} & \textbf{Keterangan dan Hasil Profil Awal} \\ \hline
\endhead
1 & ``No'' & \textit{integer} & 15 & Nomor urut baris; terisi pada 10.053 baris; bukan pengenal usaha \\ \hline
2 & ``Nama Pengusaha'' & \textit{string} & Dewi Br. Contoh & Data pribadi; tidak tersedia atribut nama usaha. Terdapat 378 \textit{record} yang namanya sama dengan \textit{record} lain setelah penyeragaman huruf; ini indikasi awal, belum tentu duplikat. Tidak diteruskan ke LLM \\ \hline
3 & ``NIK'' & \textit{char}(16) & 1234567890123456 & Data pribadi; 9.475 nilai (94,2\%) berformat 16 digit; tidak ada nilai ganda. Pengenal utama deteksi duplikasi; tidak diteruskan ke LLM \\ \hline
4 & ``Alamat'' & \textit{string} & Dusun II, Jl. Contoh & Format tidak baku; terisi pada 10.047 \textit{record}; tidak diteruskan ke LLM \\ \hline
5 & ``Desa'' & \textit{string} & Lumban Contoh & Terdapat sekitar 1.300 penulisan berbeda sehingga perlu dinormalisasi ke kode wilayah dan dicocokkan dengan kecamatan \\ \hline
6 & ``Kecamatan'' & \textit{string} & Balige & Terdapat 18 penulisan untuk 16 kecamatan (misalnya perbedaan huruf kapital); dasar rekapitulasi dan stratifikasi sampel \\ \hline
7 & ``Jenis Usaha'' & \textit{string} (teks bebas) & jualan mi gomak & Terisi pada 10.036 \textit{record} (99,8\%) dengan 2.847 isian berbeda setelah penyeragaman huruf; rata-rata 2,1 kata per isian; deskripsi pendek, informal, dan memuat istilah lokal; satu-satunya atribut bidang usaha; target standardisasi ke KBLI 2025 \\ \hline
8 & ``NIB'' & \textit{char}(13) & 1234567890123 & Hanya 11 \textit{record} (0,1\%) memuat 13 digit; 3.111 \textit{record} berisi penanda kosong (``-'', ``---'', ``\_'', ``?''); sisanya kosong. NIB tidak diwajibkan untuk pendataan sehingga tidak dihitung sebagai kekurangan kelengkapan \\ \hline
9 & ``Jumlah Modal Usaha'' & \textit{integer} (rupiah) & 2500000 & Terisi pada 9.909 \textit{record} (98,6\%); 14,4\% dari nilai terisi memuat teks atau format campuran (misalnya ``2,5 jt'') sehingga perlu dinormalisasi; dasar klasifikasi skala usaha. Atribut ini modal, bukan omzet \\ \hline
10 & ``No. Telpon'' & \textit{string} & 081200000000 & Data pribadi; bertipe \textit{string} agar angka nol di depan tidak hilang; format tidak seragam; tidak diteruskan ke LLM \\ \hline
11 & ``Keterangan'' & \textit{string} (teks bebas) & aktif & Terisi pada 421 \textit{record} (4,2\%), dan 231 di antaranya hanya berisi ``-''; isi yang bermakna bervariasi (status usaha, frekuensi berjualan, catatan wilayah); tidak dipakai sebagai masukan eksperimen \\ \hline
\end{longtable}
}

LKjIP 2025 mencatat 10.032 UMKM \cite[hlm.~50]{lkjip2025toba}, sedangkan berkas yang tersedia memuat 10.053 \textit{record} dan keterangan staf menyebut sekitar 12.000 unit pada pendataan 2020 (\lampstaf). Ketiga angka berasal dari sumber, periode, dan cara hitung yang berbeda sehingga menjadi indikasi awal persoalan konsistensi antarsumber. Angka LKjIP digunakan sebagai rujukan populasi untuk perencanaan sampel, sedangkan keterangan staf hanya memberikan konteks historis. Kerangka sampel, pengodean KBLI, dan pembanding kualitas data menggunakan rekaman yang dapat ditelusuri ke kiriman desa tahun 2026. Pada tahap \textit{data profiling}, tim memeriksa asal kiriman, menghitung jumlah rekaman per desa, dan mengidentifikasi duplikasi.

\subsection{Data Hasil Formulir Pendataan Baru}
\label{sec:data-formulir}

Data uji coba lapangan dihimpun melalui \textit{prototype} selama uji coba tiga minggu (Subbab~\ref{sec:evalpengguna}). Formulir pendataan baru mempertahankan atribut identitas dan wilayah pada Tabel~\ref{tab:dataset}, mengganti isian wilayah teks bebas dengan kode wilayah pilihan, dan menambahkan atribut yang diperlukan untuk rekapitulasi dan pemeringkatan sederhana sebagaimana Tabel~\ref{tab:formulir}. Atribut yang dihasilkan sistem dari proses standardisasi disajikan pada Tabel~\ref{tab:atribut-sistem}. Rancangan kamus data yang lengkap disusun pada tahap perancangan dan dilaporkan pada Laporan Tugas Akhir.

{\footnotesize
\setlength{\tabcolsep}{4pt}
\renewcommand{\arraystretch}{1.15}
\begin{longtable}{|>{\centering\arraybackslash}p{0.04\textwidth}|L{0.19\textwidth}|L{0.12\textwidth}|L{0.21\textwidth}|L{0.32\textwidth}|}
\caption{Atribut Formulir Pendataan Baru (Rancangan Awal)}
\label{tab:formulir}\\
\hline
\textbf{No.} & \textbf{Atribut} & \textbf{Tipe Data} & \textbf{Contoh (fiktif)} & \textbf{Keterangan} \\ \hline
\endfirsthead
\hline
\textbf{No.} & \textbf{Atribut} & \textbf{Tipe Data} & \textbf{Contoh (fiktif)} & \textbf{Keterangan} \\ \hline
\endhead
1 & Kode wilayah desa & \textit{char}(10) & 1212010001 & Dipilih dari daftar referensi wilayah; menggantikan isian teks bebas desa dan kecamatan \\ \hline
2 & Nama pengusaha & \textit{string} & Dewi Br. Contoh & Data pribadi; tidak diteruskan ke LLM \\ \hline
3 & NIK & \textit{char}(16) & 1234567890123456 & Data pribadi; divalidasi 16 digit; dasar peringatan entri ganda \\ \hline
4 & Nomor telepon & \textit{string} & 081200000000 & Data pribadi; opsional \\ \hline
5 & Nama usaha & \textit{string} & Kedai Contoh Rasa & Atribut baru; melengkapi nama pengusaha dan menjadi dasar deteksi duplikasi bersama alamat \\ \hline
6 & Jenis usaha & \textit{string} (teks bebas) & jualan mi gomak & Teks bebas singkat; target standardisasi \\ \hline
7 & Deskripsi produk utama & \textit{string} (teks bebas) & mi gomak dengan saus andaliman & Atribut baru; masukan tambahan bagi standardisasi KBLI \\ \hline
8 & Lama usaha (tahun) & \textit{integer} & 3 & Dasar penentuan wirausaha pemula \\ \hline
9 & Modal usaha (rupiah) & \textit{integer} & 2500000 & Satuan seragam rupiah penuh; dasar klasifikasi skala usaha \\ \hline
10 & Kisaran hasil penjualan per tahun & \textit{string} (kategori) & Rp50--100 juta & Atribut baru; dihimpun sebagai kisaran, bukan angka pasti; dasar klasifikasi skala usaha \\ \hline
11 & Memiliki NIB & \textit{boolean} & true & Status legalitas \\ \hline
12 & Nomor NIB & \textit{char}(13) & 1234567890123 & Opsional; divalidasi 13 digit \\ \hline
13 & Memiliki P-IRT & \textit{boolean} & false & Status legalitas \\ \hline
14 & Memiliki sertifikat halal & \textit{boolean} & false & Status legalitas \\ \hline
15 & Saluran pemasaran & \textit{string} (pilihan ganda) & warung; WhatsApp & Daftar pilihan tertutup \\ \hline
16 & Riwayat pembinaan & \textit{string} (pilihan ganda) & pelatihan digitalisasi & Mencegah pembinaan ganda \\ \hline
17 & Status operasional & \textit{string} (kategori) & aktif & Pilihan: aktif, berhenti, atau pindah \\ \hline
18 & Tanggal verifikasi & \textit{date} & 2027-03-15 & Dicatat otomatis pada setiap pembaruan; dasar pengukuran kemutakhiran \\ \hline
\end{longtable}
}

{\footnotesize
\setlength{\tabcolsep}{4pt}
\renewcommand{\arraystretch}{1.15}
\begin{longtable}{|>{\centering\arraybackslash}p{0.04\textwidth}|L{0.19\textwidth}|L{0.12\textwidth}|L{0.21\textwidth}|L{0.32\textwidth}|}
\caption{Atribut yang Dihasilkan Sistem}
\label{tab:atribut-sistem}\\
\hline
\textbf{No.} & \textbf{Atribut} & \textbf{Tipe Data} & \textbf{Contoh (fiktif)} & \textbf{Keterangan} \\ \hline
\endfirsthead
\hline
\textbf{No.} & \textbf{Atribut} & \textbf{Tipe Data} & \textbf{Contoh (fiktif)} & \textbf{Keterangan} \\ \hline
\endhead
1 & Kode KBLI & \textit{char}(5) & 56102 & Kode hasil standardisasi bila dapat ditentukan sampai tingkat kelompok; \textit{null} bila tidak \\ \hline
2 & Golongan pokok & \textit{char}(2) & 56 & Tingkat evaluasi utama \\ \hline
3 & Status standardisasi & \textit{string} & TERKODE & TERKODE atau \texttt{PERLU\_VERIFIKASI} \\ \hline
4 & Skor konsistensi & \textit{float} & 0.8 & Proporsi dari lima sampel berulang yang sama dengan keluaran utama \\ \hline
5 & Alasan singkat & \textit{string} & mi gomak termasuk penyediaan makanan & Jejak alasan bagi verifikator; maksimal dua kalimat \\ \hline
6 & Skala usaha & \textit{string} & mikro & Hasil aturan deterministik menurut PP Nomor 7 Tahun 2021 \\ \hline
\end{longtable}
}

\subsection{Data Acuan, Kamus KBLI 2025, dan Referensi Wilayah}
\label{sec:data-acuan}

\textit{Gold standard} terdiri atas 500 \textit{record} (100 pengembangan dan 400 uji) yang diambil dari data kiriman desa 2026 dan dikodekan manusia ke KBLI 2025; pemilihan, ukuran, dan penyusunannya dijelaskan pada Subbab~\ref{sec:eksperimen}. Setiap \textit{record} acuan hanya memuat pengenal pseudonim, teks jenis usaha, kecamatan, golongan pokok dan kode KBLI hasil pengodean, serta status pengodean; atribut identitas tidak disertakan. Kamus KBLI 2025 bersumber dari lampiran Peraturan BPS Nomor 7 Tahun 2025 \cite{bps2025perban7} dan menjadi basis pengetahuan \textit{retrieval}. Referensi wilayah berupa daftar kode wilayah administratif kecamatan dan desa di Kabupaten Toba yang ditetapkan Kementerian Dalam Negeri dan dipakai untuk menormalkan serta memvalidasi atribut wilayah. Berkas dinas memuat data pribadi sekitar sepuluh ribu pelaku usaha sehingga penanganannya mengikuti Subbab~\ref{sec:etika}.

\section{Dekomposisi Persoalan}
\label{sec:dekomposisi}

Persoalan utama, yaitu tidak tersedianya data UMKM yang andal untuk mendukung keputusan pembinaan, didekomposisi menjadi empat sub-persoalan yang dikaitkan dengan pertanyaan penelitian sebagaimana Tabel~\ref{tab:dekomposisi}.

{\footnotesize
\setlength{\tabcolsep}{4pt}
\renewcommand{\arraystretch}{1.15}
\begin{longtable}{|>{\centering\arraybackslash}p{0.06\textwidth}|L{0.30\textwidth}|L{0.30\textwidth}|L{0.24\textwidth}|}
\caption{Dekomposisi Persoalan Penelitian}
\label{tab:dekomposisi}\\
\hline
\textbf{Kode} & \textbf{Sub-persoalan} & \textbf{Pendekatan Solusi} & \textbf{Keterkaitan} \\ \hline
\endfirsthead
\hline
\textbf{Kode} & \textbf{Sub-persoalan} & \textbf{Pendekatan Solusi} & \textbf{Keterkaitan} \\ \hline
\endhead
SP1 & Proses penghimpunan data tidak memungkinkan pemutakhiran berkala dan pemantauan respons desa & Pendataan partisipatif melalui formulir terpandu dan kontrol akses per desa & RQ2 \\ \hline
SP2 & Deskripsi jenis usaha tidak terstandardisasi sehingga tidak dapat direkapitulasi & Standardisasi KBLI menggunakan \textit{retrieval-grounded} LLM dengan verifikasi manusia & RQ1, RQ2 \\ \hline
SP3 & Data tidak memenuhi dimensi kualitas yang diperlukan untuk keputusan (kelengkapan, akurasi, konsistensi, kemutakhiran, keunikan, dan validitas) & Validasi aturan domain, deteksi duplikasi, dan pencatatan waktu verifikasi & RQ2 \\ \hline
SP4 & Seleksi sasaran pembinaan dilakukan tanpa dukungan data terstruktur & \textit{Dashboard} rekapitulasi dan pemeringkatan sederhana (SAW) sebagai fitur pendukung keputusan yang tidak diuji & Di luar RQ1--RQ3 \\ \hline
\end{longtable}
}

\section{Kebutuhan Awal Artefak}
\label{sec:kebutuhan}

Berdasarkan Tabel~\ref{tab:temuan} dan Tabel~\ref{tab:dekomposisi}, kebutuhan fungsional (KF) dan nonfungsional (KNF) awal artefak dirumuskan pada Tabel~\ref{tab:kebutuhan}. Kebutuhan ini akan divalidasi kembali bersama dinas pada tahap penetapan tujuan solusi (Subbab~\ref{sec:tahapan}).

{\footnotesize
\setlength{\tabcolsep}{4pt}
\renewcommand{\arraystretch}{1.15}
\begin{longtable}{|>{\centering\arraybackslash}p{0.07\textwidth}|L{0.62\textwidth}|L{0.20\textwidth}|}
\caption{Kebutuhan Fungsional dan Nonfungsional Awal}
\label{tab:kebutuhan}\\
\hline
\textbf{Kode} & \textbf{Kebutuhan} & \textbf{Pengguna} \\ \hline
\endfirsthead
\hline
\textbf{Kode} & \textbf{Kebutuhan} & \textbf{Pengguna} \\ \hline
\endhead
KF-01 & Sistem menyediakan formulir terpandu untuk menambah, memperbarui, dan menandai status operasional UMKM per desa, serta mencatat tanggal verifikasi setiap pembaruan & Operator desa \\ \hline
KF-02 & Sistem memvalidasi format isian (NIB, modal, wilayah) dan memberikan peringatan entri ganda pada saat pengisian & Operator desa \\ \hline
KF-03 & Sistem mengimpor data lembar kerja kiriman desa ke dalam basis data dengan pemetaan atribut & Staf dinas \\ \hline
KF-04 & Sistem menstandardisasi deskripsi jenis usaha ke kode KBLI 2025 beserta skor konsistensi dan alasan singkat & Sistem \\ \hline
KF-05 & Sistem menyediakan antrean verifikasi untuk keluaran standardisasi yang berskor konsistensi rendah atau yang ditandai perlu verifikasi & Staf dinas \\ \hline
KF-06 & Sistem mengklasifikasikan skala usaha berdasarkan kriteria PP Nomor 7 Tahun 2021 & Sistem \\ \hline
KF-07 & Sistem menampilkan rekapitulasi per kecamatan, desa, klaster sektor, KBLI, dan skala usaha serta status pengiriman data per desa & Staf dinas, pimpinan \\ \hline
KF-08 & Sistem menampilkan pemeringkatan kandidat penerima pembinaan dengan SAW berdasarkan bobot kriteria yang ditetapkan dinas dan batasan kuota per kecamatan (fitur pendukung keputusan yang tidak diuji) & Staf dinas, pimpinan \\ \hline
KF-09 & Sistem mengekspor data dan rekapitulasi dalam format baku untuk keperluan pelaporan dan permintaan instansi lain & Staf dinas \\ \hline
KNF-01 & Kontrol akses berbasis peran; operator desa hanya mengakses data desanya & Seluruh pengguna \\ \hline
KNF-02 & Atribut data pribadi dipisahkan sebelum pemrosesan LLM dan tidak dikirim ke layanan LLM eksternal; pola nama orang dan nomor telepon pada teks bebas disamarkan sebelum dikirim; pemrosesan mematuhi UU Pelindungan Data Pribadi \cite{uupdp2022} & Sistem \\ \hline
KNF-03 & Antarmuka dapat digunakan pada perangkat seluler dan jaringan berkecepatan rendah & Operator desa \\ \hline
KNF-04 & Seluruh perubahan data dan keputusan verifikasi tercatat dalam \textit{audit trail} & Staf dinas \\ \hline
KNF-05 & Versi model, \textit{prompt}, dan kamus KBLI terdokumentasi untuk keterulangan eksperimen & Peneliti \\ \hline
KNF-06 & Latensi dan jumlah \textit{token} setiap pemanggilan LLM dicatat pada \textit{log} sebagai dasar pengukuran fase \textit{benchmarking} (Subbab~\ref{sec:benchmark}) & Peneliti \\ \hline
\end{longtable}
}

\section{Pemilihan Pendekatan}
\label{sec:analisis-metode}

\textbf{Pemilihan \textit{retrieval-grounded} LLM untuk standardisasi.} Pendekatan pembelajaran mesin terawasi memerlukan data latih berlabel dalam jumlah besar untuk banyak kelas KBLI dan menunjukkan penurunan kinerja yang tajam pada data riil \cite{amnur2025kbli}. LLM dapat dipakai pada tugas klasifikasi secara \textit{zero-shot} tanpa data latih berlabel dalam jumlah besar \cite{gilardi2023chatgpt,ziems2024css}, tetapi kinerjanya belum melampaui model terlatih khusus terbaik \cite{ziems2024css} dan LLM berisiko menghasilkan kode yang tidak sah apabila diminta menghasilkan kode dari ingatan parametriknya \cite{ji2023hallucination,dahl2024legal}. Pembatasan pilihan LLM pada kandidat hasil \textit{retrieval} dari kamus KBLI resmi dirancang untuk menekan risiko tersebut \cite{gao2023ragsurvey,achananuparp2025occupation}. Karena efektivitas rancangan ini belum diketahui untuk data UMKM berbahasa Indonesia, penelitian ini mengujinya secara eksperimental terhadap kondisi pembanding.

\textbf{Pemilihan dua jenis model.} Model \textit{proprietary} umumnya diakses melalui API sehingga pengiriman data sensitif ke pihak ketiga menjadi kendala \cite{narayan2022wrangle}, sedangkan model \textit{open-weight} dapat dijalankan pada infrastruktur lokal dan karenanya sesuai dengan keterbatasan anggaran serta persyaratan pelindungan data pemerintah daerah. Perbandingan keduanya memberikan informasi mengenai pertukaran (\textit{trade-off}) antara kinerja, biaya, dan kedaulatan data.

\textbf{Kriteria pemilihan model.} Tabel~\ref{tab:benchmark} merinci ukuran yang digunakan untuk membandingkan kandidat pada fase \textit{benchmarking} (Subbab~\ref{sec:benchmark}). Versi setiap kandidat, hasil pengukuran dan alasan pemilihan model utama serta model replikasi dicatat pada \textit{log} eksperimen.

{\footnotesize
\setlength{\tabcolsep}{4pt}
\renewcommand{\arraystretch}{1.15}
\begin{longtable}{|L{0.22\textwidth}|L{0.44\textwidth}|L{0.24\textwidth}|}
\caption{Kriteria Benchmarking LLM}
\label{tab:benchmark}\\
\hline
\textbf{Kriteria} & \textbf{Cara Ukur} & \textbf{Peran dalam Keputusan} \\ \hline
\endfirsthead
\hline
\textbf{Kriteria} & \textbf{Cara Ukur} & \textbf{Peran dalam Keputusan} \\ \hline
\endhead
Akurasi & Acc@1 golongan pokok pada kondisi P1 di 100 \textit{record} pengembangan, beserta interval kepercayaan 95\% & Kriteria utama \\ \hline
Kepatuhan format keluaran & Proporsi keluaran berupa JSON valid dengan kode yang terdapat pada daftar kandidat & Syarat kelayakan \\ \hline
Biaya & Biaya per seribu \textit{record}: token masukan dan keluaran dikalikan tarif untuk model \textit{proprietary}; perkiraan biaya perangkat keras dan listrik untuk model \textit{open-weight} & Penentu setelah akurasi; syarat kelayakan (anggaran) \\ \hline
Latensi dan \textit{throughput} & Median dan persentil ke-95 waktu per \textit{record}, serta \textit{record} per menit, pada perangkat dan jaringan yang sama & Penentu setelah biaya; syarat kelayakan (ambang) \\ \hline
Kendala operasional & Panjang konteks yang memuat $k$ kandidat, lisensi yang mengizinkan penggunaan penelitian, dan kesesuaian memori GPU untuk model \textit{open-weight} & Syarat kelayakan \\ \hline
Kinerja berbahasa Indonesia & Hasil \textit{benchmark} berbahasa Indonesia yang dilaporkan publik \cite{cahyawijaya2024cendol} & Informasi pendukung, bukan penentu \\ \hline
\end{longtable}
}

Untuk tahap \textit{retrieval}, dipilih model \textit{embedding} multibahasa yang mendukung bahasa Indonesia dengan \textit{Recall@k} tertinggi pada data pengembangan. Jumlah contoh \textit{few-shot} pada kondisi P2b ditetapkan pada data pengembangan dengan menguji tiga, lima, dan delapan contoh, lalu memilih jumlah terkecil yang tidak menurunkan Acc@1 dibandingkan jumlah yang lebih besar, karena penambahan contoh memperpanjang \textit{prompt} dan menaikkan biaya.

\textbf{Pemilihan pemeringkatan SAW.} Pendekatan pendukung keputusan pembinaan UMKM yang menggabungkan klasterisasi dan AHP-TOPSIS bekerja pada data atribut yang telah tersedia tanpa membahas standardisasi data mentah \cite{khotimah2024hybrid}. Karena penelitian ini berfokus pada penghimpunan dan standardisasi data, segmentasi dan penurunan bobot dengan AHP tidak termasuk cakupan (Subbab~\ref{sec:ruang-lingkup}). Pemeringkatan dibatasi pada SAW (Persamaan~\eqref{eq:saw}) dengan bobot yang ditetapkan langsung oleh Kepala Dinas atau Kepala Bidang karena transparan dan mudah dijelaskan kepada pimpinan dinas.

\textbf{Pemilihan pengguna aparatur, bukan pelaku UMKM.} Hasil wawancara menunjukkan bahwa alat digital mandiri bagi pelaku UMKM cenderung tidak digunakan (\lampstaf), sedangkan keberhasilan penerapan kecerdasan buatan generatif pada UMKM bergantung pada kompetensi dan kesiapan organisasi penggunanya \cite{rajaram2024genai}. Oleh karena itu, artefak ditujukan untuk memperkuat kapasitas aparatur sebagai pendamping, sesuai prinsip augmentasi \cite{raisch2021automation}.

\section{Kerangka Berpikir}
\label{sec:kerangka}

Hasil kajian awal dirangkum dalam kerangka berpikir pada Gambar~\ref{fig:kerangka}, yang menghubungkan persoalan, landasan teori, artefak yang diusulkan, evaluasi, dan luaran penelitian.

\begin{figure}[H]
  \centering
  \begin{tikzpicture}[node distance=0.5cm]
  \node (p) [boxw] {\textbf{Persoalan:} pendataan UMKM terbatas (indikator usaha mikro dan kecil 0,49\% dari target 1,80\%); pendataan berjenjang tanpa pemutakhiran berkala; jenis usaha berupa teks bebas; lima personel teknis; sumber daya pembinaan terbatas};
  \node (g) [boxw, below=of p] {\textbf{Celah penelitian:} keandalan \textit{retrieval-grounded} LLM pada data administratif berbahasa Indonesia belum teruji; belum ada bukti yang menghubungkan penghimpunan data, standardisasi, dan keputusan pembinaan pada satu artefak yang dievaluasi; dampak komponen LLM terhadap kualitas data dan kinerja tugas pengguna belum diukur berdampingan};
  \node (t2) [box, minimum height=1.5cm, below=0.6cm of g] {LLM dan RAG \cite{lotfipoor2026llm,gao2023ragsurvey}; KBLI \cite{bps2025perban7}};
  \node (t1) [box, minimum height=1.5cm, left=0.4cm of t2] {DSR \cite{hevner2004design,peffers2007dsrm}; kualitas data \cite{iso25012,govuk2021dqdimensions}};
  \node (t3) [box, minimum height=1.5cm, right=0.4cm of t2] {SUS dan NASA-TLX \cite{brooke1996sus,hernandez2022nasatlx}};
  \node (a) [boxw, below=2.2cm of g] {\textbf{Artefak:} pendataan partisipatif desa + standardisasi KBLI 2025 berbasis \textit{retrieval-grounded} LLM + validasi kualitas data + \textit{dashboard} rekapitulasi dan pemeringkatan sederhana};
  \node (e) [boxw, below=of a] {\textbf{Evaluasi:} eksperimen klasifikasi (RQ1); kualitas data sebelum--sesudah pada enam dimensi (RQ2); evaluasi \textit{within-subject} kinerja tugas dan beban kerja (RQ3)};
  \node (l) [boxw, below=of e] {\textbf{Luaran:} bukti evaluasi empiris, \textit{prototype} sistem, \textit{gold standard} KBLI beranotasi, laporan Tugas Akhir, dan manuskrip JCIS};
  \draw[arrow] (p) -- (g);
  \draw[arrow] (g) -- (t1); \draw[arrow] (g) -- (t2); \draw[arrow] (g) -- (t3);
  \draw[arrow] (t1) -- (a); \draw[arrow] (t2) -- (a); \draw[arrow] (t3) -- (a);
  \draw[arrow] (a) -- (e); \draw[arrow] (e) -- (l);
  \end{tikzpicture}
  \caption{Kerangka Berpikir Penelitian}
  \label{fig:kerangka}
\end{figure}

\section{Rancangan Artefak}
\label{sec:artefak}

\subsection{Arsitektur Sistem}

Artefak dirancang sebagai aplikasi \textit{web} dengan empat lapisan, sebagaimana Gambar~\ref{fig:arsitektur}. Lapisan antarmuka menyediakan formulir pendataan bagi operator desa serta \textit{dashboard} dan antrean verifikasi bagi staf dinas. Lapisan aplikasi menangani autentikasi, kontrol akses, validasi, dan alur verifikasi. Lapisan layanan analitis terdiri atas layanan standardisasi berbasis LLM dan layanan rekapitulasi serta pemeringkatan sederhana. Lapisan data menyimpan data UMKM, kamus KBLI 2025 beserta indeks vektornya, referensi wilayah, dan \textit{audit trail}.

\begin{figure}[H]
  \centering
  \begin{tikzpicture}[node distance=0.45cm and 0.35cm]
  \node (u1) [box, text width=4.8cm, minimum height=1.3cm] {Formulir pendataan terpandu (operator desa)};
  \node (u2) [box, text width=4.8cm, minimum height=1.3cm, right=0.8cm of u1] {\textit{Dashboard}, antrean verifikasi, dan pemeringkatan sederhana (staf dinas)};
  \path (u1.south) -- (u2.south) coordinate[midway] (mu);
  \node (ap) [boxw, below=1.0cm of mu] {\textbf{Lapisan aplikasi:} autentikasi, kontrol akses berbasis peran, validasi aturan domain, deteksi duplikasi, alur verifikasi, \textit{audit trail}, ekspor};
  \node (l1) [box, text width=4.8cm, minimum height=1.3cm, below=1.0cm of ap.south, xshift=-2.8cm] {Layanan standardisasi KBLI (\textit{retriever} hibrida + LLM)};
  \node (l2) [box, text width=4.8cm, minimum height=1.3cm, below=1.0cm of ap.south, xshift=2.8cm] {Layanan analitis (klasifikasi skala usaha, rekapitulasi, skor prioritas)};
  \path (l1.south) -- (l2.south) coordinate[midway] (ml);
  \node (db) [boxw, below=1.0cm of ml] {\textbf{Lapisan data:} basis data UMKM; kamus KBLI 2025 dan indeks vektor; referensi wilayah; bobot kriteria; \textit{audit trail}};
  \draw[arrow, <->] (u1) -- (ap); \draw[arrow, <->] (u2) -- (ap);
  \draw[arrow, <->] (ap) -- (l1); \draw[arrow, <->] (ap) -- (l2);
  \draw[arrow, <->] (l1) -- (db); \draw[arrow, <->] (l2) -- (db);
  \end{tikzpicture}
  \caption{Arsitektur Artefak}
  \label{fig:arsitektur}
\end{figure}

\subsection{\textit{Pipeline} Standardisasi KBLI}

\textit{Pipeline} standardisasi terdiri atas lima langkah. Pertama, prapemrosesan menormalkan ejaan, singkatan umum, dan format penulisan, serta memisahkan atribut data pribadi sehingga hanya teks jenis usaha, deskripsi produk (hanya tersedia pada data hasil formulir baru), dan nama kecamatan yang diteruskan ke LLM. Prapemrosesan ini identik untuk seluruh kondisi eksperimen; glosarium istilah lokal tidak diterapkan pada tahap ini, melainkan disertakan sebagai konteks pada \textit{prompt} kondisi P2a (Subbab~\ref{sec:eksperimen}). Kedua, \textit{retriever} hibrida yang menggabungkan skor leksikal BM25 dan kemiripan vektor \textit{embedding} multibahasa mengambil $k$ kandidat kelompok KBLI beserta judul dan uraiannya; nilai $k$ dipilih dari $\{5, 10, 20\}$ dengan aturan keputusan berupa nilai $k$ terkecil yang mencapai \textit{Recall@k} sekurang-kurangnya 0,90 pada tingkat golongan pokok pada data pengembangan, sepanjang panjang masukan tidak melampaui batas konteks model yang dipakai. Kisaran tersebut dipilih karena studi pemetaan deskripsi pekerjaan ke taksonomi resmi memakai jumlah kandidat pada orde yang sama \cite{achananuparp2025occupation}, sedangkan nilai $k$ yang lebih besar menambah panjang \textit{prompt} dan biaya tanpa jaminan peningkatan cakupan. Penggabungan skor BM25 dan skor kemiripan vektor dilakukan dengan \textit{reciprocal rank fusion} \cite{cormack2009rrf}, yang dipilih karena hanya memerlukan peringkat dan tidak menuntut penyetaraan skala antara kedua skor; bobot linear yang ditala pada data pengembangan dipakai sebagai pembanding, dan pilihan yang dipakai pada data uji ditetapkan berdasarkan \textit{Recall@k} pada data pengembangan. Ketiga, LLM menerima deskripsi usaha dan daftar kandidat, kemudian mengembalikan keluaran terstruktur sebagaimana Source Code~\ref{lst:prompt}; bidang \texttt{keyakinan} (\textit{confidence score} yang dinyatakan sendiri oleh model) dicatat tetapi tidak dipakai sebagai sinyal pemilah karena nilai tersebut cenderung berlebihan \cite{xiong2024calibration}. Keempat, validator memeriksa bahwa kode keluaran terdapat dalam kandidat dan kamus KBLI. Kelima, \textit{record} berstatus \texttt{PERLU\_VERIFIKASI}, atau yang skor konsistensinya (Subbab~\ref{sec:eksperimen}) berada di bawah \textit{threshold} yang ditetapkan pada data pengembangan, diteruskan ke antrean verifikasi staf dinas. Pemetaan dari kode KBLI ke klaster sektor pembinaan dinas, misalnya tenun dan kriya, olahan pangan, perdagangan, dan jasa, dilakukan menggunakan tabel pemetaan deterministik yang disepakati bersama dinas.

\begin{lstlisting}[language={}, caption={Rancangan Struktur \textit{Prompt} dan Keluaran Standardisasi KBLI}, label={lst:prompt}]
[SISTEM]
Anda adalah anotator Klasifikasi Baku Lapangan Usaha Indonesia (KBLI) 2025.
Pilih SATU kode dari DAFTAR KANDIDAT yang paling sesuai dengan kegiatan
usaha utama. Jangan membuat kode di luar daftar. Jika deskripsi tidak
cukup untuk menentukan kode, kembalikan status "PERLU_VERIFIKASI".

[GLOSARIUM LOKAL]
mi gomak = mi khas Batak; ulos = kain tenun tradisional Batak; ...

[DESKRIPSI USAHA]
jenis_usaha: "<teks dari formulir>"
produk_utama: "<teks dari formulir; kosong untuk data dinas lama>"
kecamatan: "<nama kecamatan>"

[DAFTAR KANDIDAT]
1. <kode> | <judul> | <uraian singkat>
...
k. <kode> | <judul> | <uraian singkat>

[FORMAT KELUARAN JSON]
{
  "kode_kbli": "<kode dari daftar atau null>",
  "status": "TERKODE | PERLU_VERIFIKASI",
  "keyakinan": "<bilangan 0.0 - 1.0>",
  "alasan": "<maksimal dua kalimat>"
}
\end{lstlisting}

\subsection{Modul Pendataan dan Kualitas Data}

Modul pendataan menyediakan formulir terpandu dengan daftar pilihan untuk atribut berdomain tertutup, isian bebas hanya untuk deskripsi usaha dan produk, validasi aturan domain pada saat pengisian, serta peringatan kemungkinan duplikasi berdasarkan kesamaan NIK serta kemiripan nama pengusaha, nama usaha, dan alamat dalam satu desa. Setiap pembaruan mencatat waktu verifikasi dan identitas operator sehingga kemutakhiran dapat diukur. Data lembar kerja kiriman desa diimpor melalui fungsi pemetaan atribut dan diproses oleh \textit{pipeline} standardisasi yang sama.

\subsection{Modul Pendukung Keputusan}

Modul pendukung keputusan terdiri atas dua fungsi. Fungsi rekapitulasi menyajikan jumlah dan sebaran UMKM per kecamatan, desa, klaster sektor, KBLI, dan skala usaha, serta status pengiriman data per desa. Fungsi pemeringkatan menghitung skor prioritas kandidat dengan SAW (Persamaan~\eqref{eq:saw}, Subbab~\ref{sec:spk}) menggunakan bobot kriteria yang ditetapkan langsung oleh Kepala Dinas atau Kepala Bidang tanpa penurunan bobot melalui AHP. Kandidat kriteria awal meliputi kesesuaian dengan sektor prioritas, status wirausaha pemula, kesiapan legalitas, kesiapan produk, dan riwayat penerimaan bantuan. Pemerataan antarwilayah diterapkan sebagai batasan kuota per kecamatan yang dapat diatur oleh pengguna. Fungsi ini diimplementasikan sebagai fitur pendukung keputusan agar artefak memiliki fungsi pendukung keputusan; keluarannya tidak diuji dan kesesuaiannya dengan pertimbangan ahli tidak dinilai.

\section{Rancangan Eksperimen Standardisasi (RQ1)}
\label{sec:eksperimen}

Subbab ini merancang eksperimen untuk menjawab RQ1. Pengujian utama membandingkan \textit{retrieval-grounded} LLM (P1) dengan LLM tanpa \textit{retrieval} (B1) dan \textit{baseline} leksikal (B0) berdasarkan akurasi \textit{top}-1 pada tingkat golongan pokok KBLI 2025. Dua analisis lain dilaporkan tanpa pengujian inferensial, yaitu ablasi eksploratif mengenai konteks bahasa lokal dan \textit{coverage--accuracy curve} sebagai bukti deskriptif bagi pembagian peran manusia dan mesin.

\textbf{Ukuran \textit{gold standard}.} \textit{Record} untuk \textit{gold standard} dipilih secara acak berstrata berdasarkan kecamatan dan panjang teks jenis usaha dari data UMKM kiriman desa tahun 2026 dan/atau basis data terkini dinas yang berasal dari kiriman tersebut. Kecamatan dipakai sebagai strata agar sebaran wilayah pada \textit{gold standard} mendekati sebaran populasi, sedangkan panjang teks dipakai karena isian yang sangat singkat diperkirakan lebih sulit dikodekan sehingga tingkat kesulitan perlu terwakili pada data pengembangan maupun data uji; keduanya ditetapkan sebagai strata setelah ketersediaan atributnya diperiksa pada tahap \textit{data profiling}. Sebelum pemilihan, data dideduplikasi berdasarkan identitas usaha, yaitu NIK yang sama atau kombinasi nama pengusaha, desa, dan teks jenis usaha yang sama, sehingga usaha yang sama tidak muncul pada data pengembangan dan data uji sekaligus. Ukuran data uji ditentukan agar akurasi keseluruhan dapat diestimasi dengan batas kesalahan $e = 0{,}05$ pada tingkat kepercayaan 95\% dengan asumsi proporsi paling konservatif $p = 0{,}5$, sebagaimana Persamaan~\eqref{eq:sampel}. Perhitungan menghasilkan sekitar 385 \textit{record} uji yang dibulatkan menjadi 400, sehingga \textit{gold standard} terdiri atas 500 \textit{record}, yaitu 100 \textit{record} pengembangan untuk perancangan \textit{prompt}, glosarium, contoh \textit{few-shot}, dan penetapan parameter, serta 400 \textit{record} uji yang tidak digunakan selama pengembangan. Rujukan jumlah populasi adalah 10.032 UMKM terdata pada LKjIP 2025 \cite[hlm.~50]{lkjip2025toba} (berkas yang tersedia memuat 10.053 \textit{record}; Subbab~\ref{sec:data}); dengan populasi tersebut koreksi populasi terhingga hanya menurunkan kebutuhan sampel menjadi sekitar 370 \textit{record}, sehingga koreksi tidak diterapkan, 400 \textit{record} dipertahankan, dan estimasi bersifat konservatif.

\begin{equation}
  n = \frac{z^2 \, p \, (1-p)}{e^2} = \frac{1{,}96^2 \times 0{,}5 \times 0{,}5}{0{,}05^2} = \frac{0{,}9604}{0{,}0025} \approx 385
  \label{eq:sampel}
\end{equation}
\persamaan{Ukuran sampel data uji}

Ukuran ini menentukan ketelitian estimasi akurasi, bukan daya uji. Karena RQ1 diuji dengan uji McNemar berpasangan, selisih akurasi terkecil yang dapat dideteksi bergantung pada proporsi pasangan diskordan, yaitu \textit{record} yang diklasifikasikan benar oleh satu kondisi dan salah oleh kondisi lain. Dengan $N = 400$, daya 0,80, dan $\alpha = 0{,}025$ (batas Holm terketat untuk dua perbandingan), selisih \textit{Acc@}1 minimum yang dapat dideteksi sekitar 5, 7, dan 8,5 poin persentase apabila pasangan diskordan masing-masing 10\%, 20\%, dan 30\% dari \textit{record}. Angka ini dihitung dari distribusi binomial eksak dengan jumlah pasangan diskordan dianggap tetap dan bersifat indikatif, karena proporsi diskordan sebenarnya baru diketahui setelah eksperimen. Ketelitian pada himpunan bagian, misalnya \textit{record} yang memuat istilah lokal, lebih rendah dan dilaporkan bersama interval kepercayaan 95\%.

\textbf{Penyusunan \textit{gold standard}.} Sebanyak 500 \textit{record} (100 pengembangan dan 400 uji) dikodekan ke KBLI 2025 menggunakan pedoman anotasi yang disusun tim. Pedoman menyediakan label ``informasi tidak memadai'' bagi isian yang terlalu umum untuk menentukan satu golongan pokok (Tabel~\ref{tab:risiko}, risiko 9). \textit{Record} berlabel tersebut dikeluarkan dari perhitungan Acc@1. Jumlah dan proporsi \textit{record} yang dikeluarkan serta jumlah \textit{record} efektif sebagai penyebut Acc@1 dilaporkan secara eksplisit, dengan himpunan evaluasi yang sama pada setiap kondisi pembanding. Karena itu, akurasi ditafsirkan untuk \textit{record} yang dapat diberi label, bukan seluruh \textit{record} awal. Pedoman diujicobakan pada 50 \textit{record} kalibrasi di luar \textit{gold standard} sebelum pengodean penuh. Pada \textit{record} kalibrasi ini anotator mencatat waktu pengodean per \textit{record}. Catatan tersebut dipakai untuk memperkirakan beban pemetaan manual seluruh data dinas (10.032 UMKM terdata menurut LKjIP 2025 atau jumlah \textit{record} aktual hasil \textit{data profiling}; Subbab~\ref{sec:data}) sebagai hasil kali jumlah \textit{record} dan waktu pengodean rata-rata per \textit{record}. Perkiraan ini bersifat indikatif karena dikerjakan oleh anotator dengan pedoman, bukan oleh staf dinas pada pekerjaan rutin. Anotator pertama mengodekan seluruh 500 \textit{record}, sedangkan anotator kedua secara independen mengodekan himpunan bagian acak berstrata sebesar sekitar 30\% \textit{gold standard}, yaitu 150 \textit{record} (30 dari data pengembangan dan 120 dari data uji), sehingga terdapat 650 pengodean. Ketidaksepakatan pada himpunan bagian tersebut diselesaikan melalui diskusi atau oleh anotator ketiga, dan kasus yang tetap ambigu dikonsultasikan dengan staf dinas. Pada \textit{record} di luar himpunan bagian, \textit{record} yang diragukan anotator pertama ditandai dan dikonsultasikan dengan cara yang sama.

\textbf{Justifikasi angka penyusunan \textit{gold standard}.} Proporsi pengodean ganda sebesar 30\% berada di atas kisaran 10--25\% yang dilaporkan lazim pada penilaian keandalan antarpengode \cite{oconnor2020intercoder}; proporsi yang lebih besar dipilih karena jumlah kelas KBLI banyak sehingga estimasi kappa kurang stabil. Dengan 150 \textit{record} berkode ganda, kesepakatan teramati 0,80, dan kesepakatan kebetulan 0,05, galat baku kappa sekitar 0,034 sehingga setengah lebar interval kepercayaan 95\% sekitar 0,07; ketelitian ini cukup untuk membedakan kesepakatan taraf substansial dari taraf sedang menurut batas yang lazim dirujuk \cite{landis1977kappa}. Data pengembangan sebanyak 100 \textit{record} ditetapkan sebagai jumlah terkecil yang masih memungkinkan penalaan nilai $k$, metode fusi, \textit{prompt}, dan \textit{threshold} tanpa mengurangi ukuran data uji yang telah dihitung, dan dipilih agar mencakup seluruh bidang usaha yang muncul pada data. Sebanyak 50 \textit{record} kalibrasi ditetapkan sebagai jumlah yang cukup untuk menguji keterpahaman pedoman anotasi dan mencatat waktu pengodean, tanpa mengurangi \textit{gold standard}. Anotator adalah anggota tim dan bukan pengode KBLI terlatih; karena itu staf dinas dilibatkan sebagai adjudikator bagi kasus yang tetap ambigu, dan keterbatasan ini dicatat sebagai ancaman validitas konstruk (Subbab~\ref{sec:validitas}).

Istilah lokal didefinisikan secara operasional sebagai kata atau frasa pada teks jenis usaha yang berasal dari bahasa atau dialek daerah, atau nama produk khas daerah, misalnya \textit{mi gomak} dan \textit{ulos}, yang tidak terdapat pada uraian KBLI 2025. Anotator pertama menandai \textit{record} yang memuat istilah lokal sebelum hasil model diketahui. Proporsi \textit{record} tersebut diukur pada data pengembangan dan dipakai untuk memperkirakan ukuran himpunan bagian pada data uji; bila himpunan bagian terlalu kecil untuk menghasilkan interval kepercayaan yang bermakna, ablasi dilaporkan secara deskriptif.

\textit{Inter-annotator agreement} dihitung dengan kappa Cohen \cite{rainio2024evaluation} pada himpunan bagian yang dikodekan kedua anotator, pada tingkat kategori dan golongan pokok, dan dilaporkan bersama interval kepercayaan 95\% karena himpunan bagian yang kecil menghasilkan estimasi yang lebar. Nilai ini dipakai sebagai acuan deskriptif tentang kualitas \textit{gold standard} dan batas atas kinerja yang wajar, bukan sebagai kriteria penerimaan. Keterbatasannya adalah bahwa 350 \textit{record} berlabel tunggal sehingga galat pengodean pada \textit{record} tersebut tidak terukur. Kappa model terhadap \textit{gold standard} tidak dinilai terhadap batas mutlak, melainkan dilaporkan relatif terhadap \textit{inter-annotator agreement} manusia pada himpunan bagian yang sama, karena kesepakatan manusia merupakan batas atas yang wajar bagi kesepakatan model. Batas Landis dan Koch \cite{landis1977kappa} dipakai hanya sebagai rujukan penamaan taraf kesepakatan, bukan sebagai kriteria penerimaan, karena penulisnya sendiri menyatakan batas tersebut bersifat arbitrer.

\textbf{Kondisi eksperimen.} Eksperimen menggunakan desain faktorial parsial pada data uji yang sama (400 \textit{record}) sebagaimana Tabel~\ref{tab:kondisi}. Setiap kondisi LLM dijalankan dengan dua model, yaitu satu model \textit{proprietary} dan satu model \textit{open-weight}, pada temperatur nol dan diulang tiga kali untuk mengukur keterulangan keluaran (\textit{output reproducibility}). Model utama untuk pengujian RQ1 adalah model \textit{open-weight}, karena dapat direproduksi dan dijalankan pada infrastruktur yang dikendalikan dinas. Model \textit{open-weight} tertentu dipilih melalui fase \textit{benchmarking} LLM pada data pengembangan (Subbab~\ref{sec:benchmark}) dan dicatat sebelum data uji dijalankan; model \textit{proprietary} dilaporkan sebagai replikasi. Seluruh kondisi memakai prapemrosesan yang identik, yaitu normalisasi ejaan, singkatan umum, dan format penulisan tanpa glosarium istilah lokal (Subbab~\ref{sec:artefak}); glosarium dan contoh \textit{few-shot} hanya hadir pada kondisi P2a dan P2b.

{\footnotesize
\setlength{\tabcolsep}{4pt}
\renewcommand{\arraystretch}{1.15}
\begin{longtable}{|>{\centering\arraybackslash}p{0.07\textwidth}|L{0.20\textwidth}|L{0.45\textwidth}|L{0.18\textwidth}|}
\caption{Kondisi Eksperimen Standardisasi KBLI}
\label{tab:kondisi}\\
\hline
\textbf{Kode} & \textbf{Kondisi} & \textbf{Deskripsi} & \textbf{Peran} \\ \hline
\endfirsthead
\hline
\textbf{Kode} & \textbf{Kondisi} & \textbf{Deskripsi} & \textbf{Peran} \\ \hline
\endhead
B0 & \textit{Baseline} leksikal & Kode KBLI dengan skor BM25 tertinggi terhadap uraian KBLI tanpa LLM & Pembanding RQ1 \\ \hline
B1 & LLM tanpa \textit{retrieval} & LLM diminta menghasilkan kode KBLI langsung dari deskripsi usaha (\textit{zero-shot}); dievaluasi pada tingkat golongan pokok & Pembanding RQ1 \\ \hline
B1$'$ & LLM dengan daftar golongan pokok di konteks & Daftar golongan pokok KBLI 2025 diberikan di dalam konteks tanpa \textit{retrieval}; dikerjakan bila waktu memungkinkan & Pembanding tambahan (opsional) \\ \hline
P1 & \textit{retrieval-grounded} LLM & LLM memilih kode dari $k$ kandidat hasil \textit{retriever} hibrida & Kondisi uji RQ1 \\ \hline
P2a & P1 dengan glosarium & P1 ditambah glosarium istilah lokal pada \textit{prompt} & Ablasi eksploratif \\ \hline
P2b & P1 dengan contoh \textit{few-shot} & P1 ditambah contoh \textit{few-shot} dari data pengembangan & Ablasi eksploratif \\ \hline
\end{longtable}
}

Kondisi B1 dievaluasi pada tingkat golongan pokok. Karena KBLI 2025 baru berlaku sejak diundangkan pada 18 Desember 2025 \cite{bps2025perban7}, taksonomi ini kemungkinan belum ada pada data latih model, sehingga selisih P1 terhadap B1 dapat sebagian mencerminkan ketiadaan akses terhadap taksonomi, dan bukan hanya efek pembatasan ruang keluaran. Karena itu, kondisi B1$'$ bila dikerjakan membantu memisahkan kedua efek tersebut, tetapi tidak menjadi syarat pengujian RQ1.

Parameter yang dipilih menggunakan data pengembangan (100 \textit{record}), yaitu nilai $k$, metode fusi skor \textit{retriever}, \textit{prompt}, glosarium, contoh \textit{few-shot}, pemilihan model \textit{open-weight}, dan \textit{threshold} konsistensi, dikunci sebelum data uji dijalankan. Penguncian dibuktikan dengan \textit{tag} versi dan \textit{checksum} berkas konfigurasi yang dicatat pada \textit{log} eksperimen beserta tanggalnya. Data uji tidak dipakai untuk memilih parameter apa pun, dan eksperimen pada data uji dimulai setelah seluruh anotasi data uji selesai.

\textbf{Metrik dan analisis.} Metrik primer adalah akurasi \textit{top}-1 (Acc@1) pada tingkat golongan pokok KBLI 2025 (Persamaan~\eqref{eq:topk}). Metrik sekunder, yang dilaporkan secara deskriptif dan bukan dasar kesimpulan utama RQ1, adalah kappa Cohen model terhadap \textit{gold standard} (Persamaan~\eqref{eq:kappa}), \textit{macro}-F1 (Persamaan~\eqref{eq:f1}), akurasi \textit{top}-3, tingkat halusinasi (Persamaan~\eqref{eq:halusinasi}), keterulangan keluaran, dan \textit{Recall@k} untuk tahap \textit{retrieval} (Persamaan~\eqref{eq:recallk}) pada tingkat kelompok dan golongan pokok. Keterulangan keluaran diukur sebagai proporsi \textit{record} yang memperoleh kode identik pada tiga ulangan bertemperatur nol. Tujuan ukuran ini adalah memeriksa apakah hasil eksperimen dapat diulang pada konfigurasi yang sama, bukan mengukur stabilitas model; pada temperatur nol keluaran hampir deterministik sehingga variasi yang terukur mencerminkan non-determinisme infrastruktur, misalnya urutan operasi pada GPU. Tiga ulangan dipakai karena sudah cukup untuk mendeteksi ketidakstabilan infrastruktur, sedangkan penambahan ulangan hanya menaikkan biaya inferensi. Ukuran ini berbeda dari skor konsistensi antarjawaban untuk analisis pembagian peran manusia dan mesin, yang memakai lima sampel pada temperatur lebih besar dari nol dan memang ditujukan untuk mengukur stabilitas keluaran model. Kinerja pada kode lima digit dilaporkan sebagai hasil eksploratif dan dibandingkan dengan temuan Amnur dkk. \cite{amnur2025kbli}; perbandingan ini tidak setara secara langsung karena data dan versi KBLI berbeda.

Pengujian RQ1 membandingkan P1 dengan B1 dan P1 dengan B0 pada model utama berdasarkan Acc@1 golongan pokok menggunakan uji McNemar eksak pada pasangan prediksi benar--salah (Persamaan~\eqref{eq:mcnemar}) \cite{rainio2024evaluation}. Tingkat signifikansi ditetapkan sebesar 0,05 dengan koreksi Holm untuk dua perbandingan \cite{menyhart2021multipletesting}. Hasil model \textit{proprietary} dilaporkan terpisah sebagai replikasi, dan perbandingan antarmodel diperlakukan sebagai analisis eksploratif tanpa arah yang ditetapkan lebih dahulu. Selisih akurasi, perbedaan kappa, tingkat halusinasi, dan \textit{macro}-F1 dinilai dengan interval kepercayaan 95\% hasil \textit{paired bootstrap} sebanyak 10.000 \textit{resample} pada tingkat \textit{record}, dengan \textit{seed} tetap yang dicatat pada \textit{log} eksperimen agar hasil dapat diulang. Uji McNemar eksak bersifat dua arah (Persamaan~\eqref{eq:mcnemar}) sehingga perbedaan ke kedua arah dapat terdeteksi; arah selisih dibaca dari jumlah pasangan diskordan, yaitu $b$ (P1 benar, pembanding salah) dan $c$ (P1 salah, pembanding benar), dan hasil dilaporkan apa adanya meskipun P1 tidak lebih baik daripada pembanding.

Analisis ablasi eksploratif membandingkan P2a dan P2b dengan P1 pada himpunan bagian \textit{record} yang memuat istilah lokal, dengan prapemrosesan yang identik, dan dilaporkan sebagai selisih akurasi beserta interval kepercayaan 95\% tanpa pengujian inferensial dan tanpa koreksi perbandingan ganda.

Untuk analisis pembagian peran manusia dan mesin, kondisi P1 pada model utama dilengkapi skor konsistensi $s_i$ setiap \textit{record}, yaitu proporsi dari lima sampel berulang pada temperatur lebih besar dari nol yang menghasilkan golongan pokok sama dengan keluaran utama. Konsistensi antarjawaban dilaporkan membantu mengurangi keyakinan berlebih (\textit{overconfidence}) dibandingkan \textit{confidence score} yang dinyatakan sendiri oleh model \cite{xiong2024calibration}. \textit{Record} dengan $s_i$ di bawah \textit{threshold} $\tau$ diteruskan ke verifikator. \textit{Threshold} $\tau$ ditetapkan pada data pengembangan sebagai nilai terkecil yang memberi akurasi selektif minimal sebesar target yang disepakati bersama dinas sebelum data uji dijalankan. Pada data uji dilaporkan \textit{coverage--accuracy curve} (Persamaan~\eqref{eq:cakupan}) dan proporsi \textit{record} yang diteruskan ke antrean verifikasi sebagai bukti deskriptif bagi pembagian peran tersebut. Pengukuran kalibrasi lebih lanjut, misalnya \textit{expected calibration error}, merupakan pekerjaan lanjutan. Analisis galat kualitatif dilakukan pada \textit{record} yang salah klasifikasi untuk mengidentifikasi pola, seperti istilah lokal, usaha hibrida, dan deskripsi yang terlalu umum.

\textbf{Kriteria penerimaan.} Hasil RQ1 dinyatakan mendukung efektivitas \textit{retrieval-grounded} LLM apabila uji dua arah pada Acc@1 tingkat golongan pokok menunjukkan perbedaan yang signifikan dan selisihnya positif (kondisi P1 lebih tinggi daripada B1 dan B0) pada model utama, dengan $\alpha = 0{,}05$ dan koreksi Holm untuk dua perbandingan; selisih akurasi dilaporkan beserta interval kepercayaan 95\%. Apabila kriteria tersebut tidak terpenuhi, tim kembali ke aktivitas perancangan (Gambar~\ref{fig:tahapan}) dan hasil tetap dilaporkan apa adanya. Kappa, tingkat halusinasi, dan metrik sekunder lain dilaporkan secara deskriptif dan tidak menjadi dasar kesimpulan utama RQ1.

\section{Rancangan Evaluasi Kualitas Data (RQ2)}
\label{sec:evalkualitas}

Evaluasi kualitas data menggunakan desain sebelum--sesudah pada desa percontohan. Kondisi pembanding untuk menguji RQ2 adalah data kiriman desa tahun 2026 dari desa yang sama, yang dihimpun melalui lembar kerja dinas, sedangkan kondisi sesudah diukur pada data yang dihimpun dan diverifikasi melalui artefak selama uji coba lapangan tiga minggu (Subbab~\ref{sec:evalpengguna}). Jumlah \textit{record} pada kedua kondisi ditentukan setelah \textit{data profiling} dilakukan.

Kualitas data diukur pada enam dimensi, yaitu kelengkapan, akurasi, konsistensi, kemutakhiran yang mengacu pada ISO/IEC 25012 \cite{iso25012}, serta keunikan dan validitas yang mengacu pada dimensi kualitas data Government Data Quality Hub \cite{govuk2021dqdimensions}. Pengukuran dilakukan pada atribut wajib yang tersedia pada kedua kondisi; himpunan atribut ini ditetapkan pada tahap \textit{data profiling} setelah format kedua berkas diperiksa, dan sekurang-kurangnya memuat nama pengusaha, jenis usaha, serta atribut wilayah yang tersedia (Tabel~\ref{tab:dataset}). Kelengkapan, validitas, keunikan, dan kemutakhiran dihitung menggunakan Persamaan~\eqref{eq:completeness}, \eqref{eq:validity}, \eqref{eq:uniqueness}, dan \eqref{eq:currentness}. Dua dimensi lainnya (persamaannya pada Subbab~\ref{sec:kualitas}) dioperasionalkan sebagai berikut. Akurasi adalah proporsi nilai atribut terisi pada sampel acak yang sesuai dengan keadaan sebenarnya berdasarkan verifikasi silang kepada perangkat desa atau pelaku usaha, sebagaimana Persamaan~\eqref{eq:akurasi}, dengan $S$ adalah jumlah nilai terisi yang diverifikasi dan $B$ adalah jumlah nilai yang sesuai. Ukuran sampel verifikasi dihitung agar selisih yang diharapkan dapat dideteksi. Untuk mendeteksi peningkatan akurasi dari 0,80 menjadi 0,92 pada $\alpha = 0{,}05$ dua arah dan daya 0,80, uji dua proporsi memerlukan sekitar 128 nilai per kondisi. Karena nilai bersarang dalam \textit{record} dan \textit{record} bersarang dalam tiga sampai empat desa, angka tersebut dikalikan faktor rancangan (\textit{design effect}) sebesar 1,5 sehingga menjadi sekitar 192 nilai per kondisi; faktor rancangan ditetapkan ulang setelah korelasi intraklas diperkirakan dari data uji coba. Skenario minimum apabila data yang tersedia terbatas adalah 50 \textit{record} per kondisi, yang pada $p \approx 0{,}9$ memberi margin kesalahan sekitar delapan poin persentase; pada skenario ini dimensi akurasi dilaporkan secara deskriptif tanpa pengujian inferensial. Hasil pada setiap skenario dilaporkan bersama interval kepercayaan 95\%.

Konsistensi adalah proporsi \textit{record} yang tidak melanggar aturan konsistensi antaratribut dan antar-\textit{record}, sebagaimana Persamaan~\eqref{eq:konsistensi}, dengan $I$ adalah jumlah \textit{record} yang melanggar sekurang-kurangnya satu aturan. Aturan ditetapkan bersama dinas pada tahap \textit{data profiling}, misalnya kesesuaian desa dengan kecamatan menurut referensi wilayah, kesesuaian skala usaha dengan modal menurut PP Nomor 7 Tahun 2021 \cite{pp7tahun2021}, serta keseragaman satuan dan format penulisan nilai atribut yang sama antar-\textit{record}.

Pengukuran dilakukan dengan skrip penilai yang terpisah dari validator artefak dan diterapkan secara identik pada kedua kondisi agar nilai sesudah tidak sekadar mencerminkan aturan validasi yang dibangun ke dalam formulir. Keberadaan kode KBLI 2025 yang valid dan terverifikasi dilaporkan sebagai indikator deskriptif tambahan pada kondisi sesudah karena atribut tersebut tidak tersedia pada kondisi sebelum.

RQ2 diuji secara inferensial pada dua dimensi, yaitu akurasi dan keunikan, menggunakan uji dua arah pada tingkat signifikansi 0,05 dengan koreksi Holm untuk dua uji. Unit analisis ditetapkan pada tingkat \textit{record}. Karena \textit{record} bersarang dalam desa sehingga asumsi independensi pada uji dua proporsi biasa tidak terpenuhi, selisih proporsi diuji dengan \textit{cluster bootstrap} yang meresampel desa sebanyak 10.000 kali dengan \textit{seed} tetap, dan model bertingkat (\textit{generalized linear mixed model}) dengan efek acak desa dipakai sebagai pemeriksaan pembanding apabila jumlah desa memungkinkan. Peningkatan dinyatakan terbukti apabila kedua dimensi menunjukkan perbedaan yang signifikan setelah koreksi dengan selisih positif (nilai sesudah lebih tinggi daripada data kiriman desa 2026). Keunikan dihitung atas identitas usaha, yaitu NIK yang sama atau kombinasi nama pengusaha, atribut wilayah, dan jenis usaha yang sama, bukan atas NIK saja, karena pengenal tunggal yang hampir selalu unik dapat menyembunyikan duplikasi usaha. Ambang keunikan kondisi pembanding yang membuat peningkatan tidak lagi dapat dideteksi ditetapkan pada tahap \textit{data profiling} sebelum data uji dijalankan; apabila keunikan pembanding sudah berada di atas ambang tersebut, dimensi keunikan dilaporkan secara deskriptif dan RQ2 diuji pada dimensi akurasi saja. Empat dimensi lainnya (kelengkapan, konsistensi, kemutakhiran, dan validitas) dilaporkan secara deskriptif disertai interval kepercayaan 95\%, karena sebagian besar nilainya ditentukan oleh rancangan formulir pendataan. Khusus kemutakhiran, format lembar kerja dinas tidak memuat atribut tanggal verifikasi data (Subbab~\ref{sec:latar-belakang}), sehingga $R_{90} = 0$ dan nilai kemutakhiran pada kondisi sebelum bernilai nol menurut definisi pengukurannya (Persamaan~\eqref{eq:currentness}); perbandingan sebelum--sesudah pada dimensi ini tidak informatif untuk diuji dan hanya dilaporkan secara deskriptif. Karena \textit{record} bersarang dalam desa dan jumlah desa hanya tiga sampai empat, hasil juga dilaporkan per desa dan ditafsirkan sebagai bukti pada studi kasus. Karena pembanding adalah proses lembar kerja secara menyeluruh, sedangkan formulir menerapkan validasi saat pengisian, peringatan entri ganda, pengendalian mutu pada titik asal data, dan antrean verifikasi sekaligus, peningkatan akurasi dan keunikan ditafsirkan sebagai efek gabungan fitur \textit{prototype}, bukan efek satu fitur secara terpisah. Selain itu, diukur dua indikator proses, yaitu jumlah \textit{record} yang dihimpun dan diverifikasi per desa selama uji coba serta waktu yang diperlukan sejak data dikirim hingga tersedia dalam rekapitulasi dinas. Cakupan populasi UMKM (proporsi usaha yang terdata) tidak termasuk ukuran penelitian ini.

\section{Rancangan Evaluasi Pengguna (RQ3)}
\label{sec:evalpengguna}

\textbf{Partisipan dan lokasi.} Desa percontohan dipilih secara purposif bersama dinas sebanyak tiga sampai empat desa dengan karakteristik berbeda, yaitu desa dengan konsentrasi usaha tenun dan kriya serta desa yang relatif jauh dari ibu kota kabupaten. Partisipan terdiri atas operator desa pada desa tersebut, sebagai kelompok analisis utama dengan target enam sampai delapan orang (satu sampai dua operator pada setiap desa), dan dua orang staf teknis dinas yang melaksanakan standardisasi dan verifikasi, yang dilaporkan secara deskriptif. Jumlah staf dinas dibatasi dua orang karena tenaga teknis yang menangani UMKM hanya lima orang dan tetap harus menangani tugas lain (\lampstaf).
\textbf{Analisis daya.} Jumlah partisipan RQ3 dibatasi oleh ketersediaan operator, bukan ditetapkan dari analisis daya a priori. Dengan tiga sampai empat desa percontohan dan satu sampai dua operator pada setiap desa, tersedia enam sampai delapan operator. Sebagai pembanding, dengan ukuran efek berpasangan $d_z = 0{,}8$, $\alpha = 0{,}05$ satu arah, dan daya 0,80, uji $t$ berpasangan memerlukan dua belas partisipan menurut distribusi $t$ nonsentral (daya 0,829 pada $n = 12$ dan 0,795 pada $n = 11$). Karena \textit{Wilcoxon signed-rank test} memiliki efisiensi asimtotik relatif sekitar 0,955 terhadap uji $t$ pada sebaran normal, kebutuhan tersebut menjadi tiga belas partisipan. Pada enam sampai delapan operator, daya uji $t$ berpasangan dengan ukuran efek yang sama hanya sekitar 0,52 ($n = 6$) sampai 0,65 ($n = 8$), dan daya \textit{Wilcoxon signed-rank test} diperkirakan sedikit lebih rendah (sekitar 0,49 sampai 0,58 menurut simulasi pada Tabel~\ref{tab:daya-wilcoxon}). Oleh karena itu RQ3 diperlakukan sebagai pengujian eksploratif dengan daya terbatas, bukan pengujian konfirmatori. Signifikansi masih dapat dicapai karena nilai $p$ satu arah terkecil adalah $1/64 \approx 0{,}016$ pada $n = 6$ dan $1/256 \approx 0{,}004$ pada $n = 8$, tetapi hanya apabila seluruh selisih berpasangan searah. Nilai $p$ eksak karenanya dilaporkan bersama ukuran efek $r$ dan interval kepercayaan, dan kesimpulan dinyatakan sebagai indikasi awal. Ukuran efek besar dipakai karena perbandingan dilakukan antara proses manual berbasis lembar kerja dan formulir terpandu yang memuat daftar pilihan serta validasi otomatis, sehingga selisih yang diharapkan tidak kecil; apabila ukuran efek yang sebenarnya lebih kecil, daya uji menurun dan hal ini dicatat sebagai keterbatasan. Apabila operator yang tersedia kurang dari enam orang, RQ3 dilaporkan secara deskriptif disertai ukuran efek dan interval kepercayaan, tanpa kesimpulan inferensial.

\textbf{Definisi konstruk kinerja tugas.} Konstruk yang diuji RQ3 adalah efisiensi entri data pada tingkat operator, bukan efisiensi proses pendataan ujung ke ujung. Tugas terkontrol memasukkan sepuluh data UMKM tidak merepresentasikan rantai pelaporan berjenjang dari desa ke kecamatan hingga dinas, sehingga kesimpulan RQ3 dibatasi pada efisiensi entri. Efisiensi proses dilaporkan sebagai ukuran pendamping tanpa pengujian inferensial, yaitu waktu sejak data dikirim operator hingga tersedia pada rekapitulasi dinas (\textit{lead time}) pada kedua kondisi, beserta jumlah langkah penyerahan pada masing-masing proses.

\textbf{Alur uji coba lapangan.} Uji coba berlangsung tiga minggu sebagaimana Tabel~\ref{tab:alur}. Durasi tiga minggu dipilih karena keterbatasan jadwal semester genap (Tabel~\ref{tab:jadwal2}) dan karena cukup untuk satu siklus pendataan per desa beserta verifikasi oleh dua orang staf dinas, sedangkan jumlah tiga sampai empat desa dipilih sebagai jumlah terbesar yang masih dapat didampingi dan dilatih langsung oleh tim pada rentang waktu tersebut, dengan satu sampai dua operator pada setiap desa sehingga tersedia enam sampai delapan operator. Sesi tugas terkontrol pada minggu pertama menghasilkan data RQ3, sedangkan pendataan nyata selama tiga minggu menghasilkan data kondisi sesudah untuk RQ2.

{\footnotesize
\setlength{\tabcolsep}{4pt}
\renewcommand{\arraystretch}{1.15}
\begin{longtable}{|>{\centering\arraybackslash}p{0.10\textwidth}|L{0.52\textwidth}|L{0.26\textwidth}|}
\caption{Alur Uji Coba Lapangan Tiga Minggu}
\label{tab:alur}\\
\hline
\textbf{Minggu} & \textbf{Kegiatan} & \textbf{Data yang Dihasilkan} \\ \hline
\endfirsthead
\hline
\textbf{Minggu} & \textbf{Kegiatan} & \textbf{Data yang Dihasilkan} \\ \hline
\endhead
1 & Orientasi dan pelatihan singkat; sesi tugas terkontrol pada dua kondisi (lembar kerja dan \textit{prototype}) dengan urutan diseimbangkan; pengisian NASA-TLX dan SUS setelah setiap kondisi & Waktu tugas, jumlah kesalahan, RTLX, SUS \\ \hline
1--3 & Pendataan nyata oleh operator desa menggunakan \textit{prototype}; dua orang staf dinas memverifikasi \textit{record} dalam antrean verifikasi (berstatus \texttt{PERLU\_VERIFIKASI} atau berskor konsistensi rendah) & \textit{Record} kondisi sesudah (RQ2), catatan penggunaan \\ \hline
3 & Verifikasi silang sampel \textit{record}; wawancara pascatugas & Akurasi data (RQ2), data kualitatif \\ \hline
\end{longtable}
}

\textbf{Jalur cadangan.} Apabila izin pelaksanaan di desa tidak terbit sebelum titik keputusan G3 (Tabel~\ref{tab:gate}), uji coba dialihkan menjadi studi terkontrol dengan staf dinas dan perwakilan desa menggunakan data kiriman desa tahun 2026.
Pada jalur cadangan ini \textit{record} kondisi sesudah berasal dari tugas terkontrol dan partisipannya staf dinas serta perwakilan desa, sehingga dampaknya pada RQ2 dan RQ3 dilaporkan secara eksplisit sebagai keterbatasan.

\textbf{Desain dan tugas.} Evaluasi menggunakan desain \textit{within-subject} dengan dua kondisi, yaitu proses berbasis lembar kerja yang berjalan saat ini dan proses menggunakan \textit{prototype}. Urutan kondisi diseimbangkan (\textit{counterbalanced}) dan materi tugas pada kedua kondisi dibuat setara tetapi tidak identik untuk mengurangi efek pembelajaran. Operator desa mengerjakan tugas memasukkan dan memutakhirkan sepuluh data UMKM. Dua orang staf dinas mengerjakan tugas mengodekan sejumlah deskripsi usaha ke KBLI 2025 dan menyusun rekapitulasi per kecamatan dan jenis usaha; pada kondisi \textit{prototype}, tugas dikerjakan melalui antrean verifikasi dan rekapitulasi otomatis.

\textbf{Ukuran.} Ukuran yang digunakan adalah waktu penyelesaian tugas dalam detik, jumlah kesalahan terhadap kunci jawaban yang disusun sebelum pengambilan data, skor \textit{raw} NASA-TLX setelah setiap kondisi (Persamaan~\ref{eq:rtlx}), dan skor SUS setelah setiap kondisi (Persamaan~\ref{eq:sus}) \cite{brooke1996sus,hernandez2022nasatlx}. Wawancara pascatugas dilakukan untuk menggali alasan di balik skor dan hambatan penggunaan.

\textbf{Integrasi data kuantitatif dan kualitatif.} Penelitian ini menggunakan desain sekuensial eksplanatori (\textit{explanatory sequential}) untuk menghubungkan pengukuran kinerja operator dengan penjelasan yang diperoleh melalui wawancara \cite{storey2025mixed}. Data kuantitatif berupa waktu penyelesaian tugas, jumlah kesalahan, skor RTLX, dan skor SUS dikumpulkan dan dianalisis lebih dahulu, kemudian wawancara pascatugas dipakai untuk menjelaskan pola yang muncul, misalnya kondisi yang skornya berlawanan dengan arah yang diharapkan. Pedoman wawancara disusun setelah hasil kuantitatif sementara tersedia dan memuat pertanyaan terbuka mengenai hambatan penggunaan, langkah yang dirasa paling lambat, dan tingkat kepercayaan terhadap keluaran standardisasi. Transkrip dianalisis dengan analisis tematik enam langkah Braun dan Clarke \cite{braun2006thematic}, yaitu pengenalan data, pengodean awal, pencarian tema, penelaahan tema, penamaan tema, dan pelaporan. Pengodean dilakukan satu anggota tim dan ditelaah anggota kedua, dan perbedaan diselesaikan melalui diskusi. Integrasi dilaporkan dalam satu tabel yang menyandingkan setiap temuan kuantitatif dengan kutipan yang menjelaskannya.

\textbf{Analisis.} Karena jumlah partisipan kecil dan sebaran selisih tidak diasumsikan normal, RQ3 diuji menggunakan \textit{Wilcoxon signed-rank test} satu arah \cite{wilcoxon1945individual,rainio2024evaluation} (Persamaan~\eqref{eq:wilcoxon}) dengan nilai $p$ eksak terhadap selisih berpasangan waktu penyelesaian tugas dan skor RTLX pada operator desa, dengan ukuran efek $r$. Uji satu arah dipakai hanya pada RQ3 karena analisis ini bersifat eksploratif dengan daya terbatas (Analisis daya di atas) dan arah yang diharapkan ditetapkan sebelum pengambilan data, yaitu formulir terpandu dengan validasi tidak diharapkan menambah waktu penyelesaian maupun beban kerja dibandingkan lembar kerja \cite{bland1994onetwosided}; pada $n = 6$ uji dua arah hanya dapat mencapai $p = 2/64 \approx 0{,}031$ dan daya ujinya lebih rendah lagi. Apabila arah yang teramati berlawanan, hasil dilaporkan apa adanya sebagai temuan yang tidak sesuai dengan arah yang diharapkan, dan nilai $p$ dua arah eksak dilaporkan pula sebagai analisis sensitivitas untuk kedua ukuran. Hasil RQ3 dinyatakan mendukung secara eksploratif apabila kedua uji menunjukkan penurunan yang signifikan pada tingkat 0,05; tanpa koreksi ganda karena kedua syarat harus terpenuhi sekaligus. Mengingat daya uji terbatas pada enam sampai delapan operator (Analisis daya di atas), hasil ini ditafsirkan sebagai indikasi awal, bukan bukti konfirmatori. Apabila operator desa kurang dari enam orang, analisis dilaporkan secara deskriptif disertai ukuran efek dan tanpa kesimpulan inferensial. Karena pembanding adalah proses lembar kerja secara menyeluruh, penurunan waktu dan beban kerja ditafsirkan sebagai efek gabungan fitur \textit{prototype}. Jumlah kesalahan dan skor SUS dilaporkan secara deskriptif dan tidak dipakai sebagai dasar kesimpulan utama RQ3. Hasil staf dinas dilaporkan secara deskriptif disertai kutipan wawancara, dan data wawancara dirangkum menurut tema yang muncul. Hasil ditafsirkan sebagai bukti efektivitas dalam konteks studi kasus, bukan generalisasi statistik.

\section{Validitas Penelitian}
\label{sec:validitas}

Matthay dan Glymour \cite{matthay2020validity} mengelompokkan ancaman terhadap validitas menjadi empat jenis, yaitu validitas internal, konstruk, eksternal, dan kesimpulan statistik. Verdecchia dkk.\ \cite{verdecchia2023validity} menekankan perlunya mengaitkan ancaman dengan keputusan dalam rancangan studi. Untuk setiap jenis validitas tersebut, penelitian ini mengidentifikasi bagian evaluasi yang rentan dan menjelaskan bagaimana mitigasinya membatasi penafsiran hasil.

\textbf{Validitas internal.} Evaluasi kualitas data (RQ2) memakai desain sebelum--sesudah pada satu kelompok, yang rentan terhadap pengaruh peristiwa di luar intervensi (\textit{history}), perubahan alamiah pada unit yang diamati (\textit{maturation}), perubahan perilaku karena merasa diamati (efek Hawthorne), dan efek kebaruan. Mitigasinya adalah memakai data kiriman desa tahun 2026 pada periode berdekatan sebagai pembanding utama; melaporkan hasil per desa; serta mencatat kejadian di luar rencana selama uji coba pada catatan lapangan dan melaporkannya bersama hasil. Efek kebaruan tidak dapat dihilangkan pada uji coba tiga minggu dan dinyatakan sebagai keterbatasan. Desain \textit{within-subject} pada RQ3 rentan terhadap efek bawaan (\textit{carryover}) meskipun urutan kondisi diseimbangkan; mitigasinya adalah jeda antarkondisi, materi tugas yang setara tetapi tidak identik, dan pelaporan pengaruh urutan sebagai analisis tambahan. Pada RQ1, kebocoran data uji dicegah melalui aturan penguncian konfigurasi (Subbab~\ref{sec:eksperimen}).

\textbf{Validitas konstruk.} Konstruk kinerja tugas dibatasi pada efisiensi entri data dan tidak mencakup efisiensi proses ujung ke ujung (Subbab~\ref{sec:evalpengguna}). Akurasi pengodean KBLI diukur terhadap \textit{gold standard} yang disusun anotator tim yang bukan pengode KBLI terlatih, sehingga galat pedoman anotasi dapat terbawa ke ukuran akurasi; mitigasinya adalah pengodean ganda pada 30\% \textit{record}, kalibrasi pedoman, dan adjudikasi oleh staf dinas untuk kasus ambigu. Kappa Cohen tidak dipakai sebagai ukuran kinerja model karena pada jumlah kelas yang banyak nilainya mendekati akurasi. Penilaian akurasi data pada RQ2 dilakukan oleh penilai yang tidak mengetahui kondisi asal \textit{record} (\textit{blinded}), dan \textit{record} dari kedua kondisi dinilai dalam urutan teracak.

\textbf{Validitas eksternal.} Penelitian dilaksanakan pada satu kabupaten dengan tiga sampai empat desa percontohan dan satu versi taksonomi, sehingga hasilnya tidak digeneralisasi ke kabupaten lain, ke taksonomi lain, maupun ke model yang tidak diuji. Keberlakuan temuan dibatasi pada batas generalisasi yang dinyatakan pada Subbab~\ref{sec:ruang-lingkup}. Untuk memperjelas batas keberlakuan, karakteristik desa percontohan, versi model, dan konfigurasi \textit{pipeline} dilaporkan lengkap.

\textbf{Validitas kesimpulan statistik.} Jumlah partisipan RQ3 kecil sehingga daya uji terbatas; mitigasinya adalah analisis daya a priori, pelaporan ukuran efek dan interval kepercayaan, serta kesiapan melaporkan RQ3 secara deskriptif apabila jumlah partisipan tidak terpenuhi. Pengujian berganda dikoreksi dengan metode Holm \cite{menyhart2021multipletesting}. Ketergantungan antar-\textit{record} dalam satu desa ditangani dengan \textit{cluster bootstrap} (Subbab~\ref{sec:evalkualitas}). Metrik sekunder dilaporkan secara deskriptif dan tidak dipakai sebagai dasar kesimpulan utama, sehingga tidak menambah risiko kesalahan jenis pertama.

\textbf{Peran peneliti.} Tim berperan sebagai pengembang artefak sekaligus evaluator, yang dapat memengaruhi penilaian. Mitigasinya adalah pemisahan peran di dalam tim, yaitu anggota yang mengimplementasikan modul pendataan tidak menjadi penilai akurasi data pada kondisi yang sama (Subbab~\ref{sec:peran}); penggunaan kunci jawaban dan pedoman penilaian yang disusun sebelum pengambilan data; serta pencatatan seluruh keputusan analisis pada \textit{log} sebelum data uji dijalankan.

\section{Pemetaan Pertanyaan Penelitian}
\label{sec:pemetaan}

Pemetaan setiap pertanyaan penelitian terhadap kegiatan, luaran, dan cara evaluasi disajikan pada Tabel~\ref{tab:pemetaan_rq}.

{\footnotesize
\setlength{\tabcolsep}{4pt}
\renewcommand{\arraystretch}{1.15}
\begin{longtable}{|>{\centering\arraybackslash}p{0.06\textwidth}|L{0.30\textwidth}|L{0.25\textwidth}|L{0.29\textwidth}|}
\caption{Pemetaan Pertanyaan Penelitian terhadap Kegiatan Penelitian}
\label{tab:pemetaan_rq}\\
\hline
\textbf{RQ} & \textbf{Kegiatan} & \textbf{Luaran} & \textbf{Cara Evaluasi} \\ \hline
\endfirsthead
\hline
\textbf{RQ} & \textbf{Kegiatan} & \textbf{Luaran} & \textbf{Cara Evaluasi} \\ \hline
\endhead
RQ1 & Penyusunan \textit{gold standard} 500 \textit{record}; perancangan \textit{retriever} dan \textit{prompt}; eksperimen B0--P2b pada model utama \textit{open-weight} dan replikasi pada model \textit{proprietary} & \textit{Gold standard} beranotasi, \textit{pipeline} standardisasi, laporan eksperimen & Uji utama: Acc@1 golongan pokok dengan uji McNemar eksak dan koreksi Holm; metrik sekunder (kappa Cohen, \textit{macro-F1}, akurasi \textit{top-3}, tingkat halusinasi, keterulangan keluaran, \textit{Recall@k}) dan ablasi P2a--P2b secara deskriptif dengan interval kepercayaan \textit{bootstrap} \\ \hline
RQ2 & \textit{Data profiling} pada data kiriman desa 2026; uji coba lapangan tiga minggu; pengukuran kualitas data kiriman desa 2026 dan data hasil artefak & Laporan kualitas data & Uji utama: akurasi dan keunikan terhadap data kiriman desa 2026 dengan uji dua proporsi dan koreksi Holm untuk dua uji; kelengkapan, konsistensi, kemutakhiran, dan validitas dilaporkan deskriptif; verifikasi silang sampel \\ \hline
RQ3 & Evaluasi \textit{within-subject} pada operator desa; staf dinas dilaporkan secara deskriptif; wawancara pascatugas & Laporan evaluasi pengguna & Uji utama: waktu tugas dan RTLX dengan uji Wilcoxon eksak (eksploratif; deskriptif bila operator kurang dari enam orang); jumlah kesalahan, SUS, dan wawancara secara deskriptif \\ \hline
\end{longtable}
}

\section{Pembagian Peran Tim}
\label{sec:peran}

Penelitian dilaksanakan oleh tiga mahasiswa dengan beban 4 SKS pada semester ganjil dan 4 SKS pada semester genap per mahasiswa. Pembagian peran dirancang agar setiap mahasiswa memiliki tanggung jawab utama atas satu pertanyaan penelitian dan satu komponen artefak, sekaligus berkontribusi pada integrasi sistem dan penulisan bersama, sebagaimana Tabel~\ref{tab:peran}.

{\footnotesize
\setlength{\tabcolsep}{4pt}
\renewcommand{\arraystretch}{1.15}
\begin{longtable}{|L{0.15\textwidth}|L{0.37\textwidth}|L{0.37\textwidth}|}
\caption{Pembagian Peran dan Beban Kerja Tim}
\label{tab:peran}\\
\hline
\textbf{Mahasiswa} & \textbf{Semester Ganjil (4 SKS)} & \textbf{Semester Genap (4 SKS)} \\ \hline
\endfirsthead
\hline
\textbf{Mahasiswa} & \textbf{Semester Ganjil (4 SKS)} & \textbf{Semester Genap (4 SKS)} \\ \hline
\endhead
M1: Analis proses bisnis dan kualitas data (penanggung jawab RQ2) & Pengurusan izin dan akses data; pemodelan proses \textit{as-is} dan \textit{to-be}; validasi kebutuhan; \textit{data profiling} dan pengukuran kualitas data kiriman desa 2026; rancangan formulir pendataan & Implementasi modul pendataan dan validasi; koordinasi uji coba lapangan dan pelatihan operator desa; evaluasi kualitas data sebelum--sesudah; verifikasi silang sampel; anotator kedua (30\% \textit{gold standard}) \\ \hline
M2: Perekayasa \textit{pipeline} LLM (penanggung jawab RQ1) & Pedoman anotasi KBLI; kalibrasi dan \textit{gold standard} pengembangan; indeks KBLI dan \textit{retriever}; eksperimen awal pada data pengembangan; \textit{benchmarking} LLM (akurasi, biaya, latensi) dan pemilihan model & Penyelesaian \textit{gold standard} uji; eksperimen B0--P2b pada model utama \textit{open-weight} dan replikasi pada model \textit{proprietary}; analisis statistik dan galat; integrasi layanan standardisasi dan antrean verifikasi \\ \hline
M3: Perekayasa analitik dan pendukung keputusan (penanggung jawab RQ3) & Perancangan arsitektur, basis data, dan antarmuka; penetapan kriteria dan bobot pemeringkatan bersama dinas; rancangan instrumen evaluasi pengguna & Implementasi \textit{dashboard} rekapitulasi dan pemeringkatan sederhana; evaluasi \textit{within-subject} dan analisis RQ3; anotator ketiga (adjudikasi) \\ \hline
Bersama & Studi literatur; wawancara; proposal; laporan TA-1 & \textit{Functional testing}; laporan Tugas Akhir; manuskrip JCIS \\ \hline
\end{longtable}
}

\section{Etika Penelitian dan Pelindungan Data}
\label{sec:etika}

Penelitian ini memproses data yang memuat informasi pribadi pelaku usaha, sehingga pelaksanaannya mengikuti prinsip UU Nomor 27 Tahun 2022 tentang Pelindungan Data Pribadi \cite{uupdp2022}. Langkah yang ditempuh meliputi: (1) penggunaan data dan pelaksanaan uji coba berdasarkan surat izin dan persetujuan tertulis dinas serta izin dari desa yang menjadi lokasi uji coba; (2) pemisahan atribut identitas dari atribut yang diproses LLM, sehingga layanan LLM eksternal hanya menerima teks jenis usaha, deskripsi produk, dan nama kecamatan; (3) penggunaan model \textit{open-weight} pada infrastruktur lokal untuk pemrosesan yang tidak dapat dipisahkan dari data pribadi; (4) penyimpanan data pada \textit{server} dengan kontrol akses dan enkripsi; (5) persetujuan setelah penjelasan (\textit{informed consent}) dari partisipan evaluasi pengguna; serta (6) publikasi \textit{gold standard} hanya dalam bentuk teranonimkan setelah memperoleh persetujuan dinas. Kepatuhan terhadap butir (2) diperiksa melalui daftar periksa yang membandingkan seluruh masukan yang dikirim ke layanan LLM selama eksperimen dan uji coba dengan atribut yang diizinkan, yaitu teks jenis usaha, deskripsi produk, dan nama kecamatan (KNF-02), serta memeriksa \textit{audit trail} untuk memastikan setiap keputusan akhir atas keluaran yang diteruskan ke antrean verifikasi dicatat atas nama verifikator manusia (KNF-04). Hasil pemeriksaan ini dilaporkan sebagai bukti pemenuhan kebutuhan nonfungsional KNF-02 dan KNF-04, bukan sebagai bukti efektivitas komponen LLM. Berkas lembar kerja dinas memuat identitas sekitar sepuluh ribu pelaku usaha (nama, NIK, alamat, dan nomor telepon). Berkas asli disimpan terenkripsi dengan akses terbatas pada tim peneliti, sedangkan salinan kerja yang dipakai untuk profil, pengambilan sampel, dan eksperimen memuat pengenal pseudonim sebagai pengganti atribut identitas. Proposal, laporan, dan manuskrip tidak memuat data pelaku usaha yang sebenarnya; seluruh contoh bersifat fiktif atau tersamarkan. Penggunaan kecerdasan buatan generatif dalam penelitian dan penulisan akan diungkapkan sesuai ketentuan jurnal sasaran.

Dalam transkrip yang disertakan pada naskah, nama orang ditulis dengan inisial. Persetujuan staf untuk diwawancarai tidak mencakup pencantuman nama pihak ketiga. Tim memeriksa kutipan, identitas penutur, dan rincian yang dapat menunjuk langsung kepada seseorang sebelum naskah dibagikan (\lampstaf{} dan \lampkadis).

Tiga langkah tambahan ditempuh. Pertama, persetujuan etik penelitian diajukan kepada unit yang berwenang di Institut Teknologi Del sebelum pengambilan data yang melibatkan manusia dimulai, yaitu sebelum sesi evaluasi pengguna dan sebelum verifikasi silang yang menghubungi pelaku usaha; apabila unit tersebut menyatakan penelitian ini tidak memerlukan telaah etik, pernyataan tersebut dilampirkan pada Laporan Tugas Akhir. Kedua, pelaku usaha yang dihubungi untuk verifikasi silang akurasi data diberi penjelasan mengenai tujuan penelitian, data yang akan dikonfirmasi, dan hak untuk menolak, lalu persetujuannya direkam secara lisan dan dicatat pada lembar verifikasi; penolakan dicatat sebagai data tidak terverifikasi dan tidak digantikan dengan dugaan. Ketiga, karena teks bebas pada atribut jenis usaha dapat memuat data pribadi, misalnya nama pemilik yang menjadi bagian nama warung, \textit{pipeline} menjalankan deteksi dan penyamaran pola nama orang serta nomor telepon pada teks tersebut sebelum teks dikirim ke layanan LLM. Langkah ini menjadi bagian daftar periksa KNF-02, dan proporsi teks yang tersamarkan dilaporkan bersama hasil pemeriksaan kepatuhan.

\section{Alat dan Bahan Penelitian}
\label{sec:alat}

Alat dan bahan penelitian disajikan pada Tabel~\ref{tab:tools}. Model dipilih melalui fase \textit{benchmarking} LLM pada semester ganjil (Subbab~\ref{sec:benchmark}); versi spesifik model dan pustaka dikunci sebelum data uji dijalankan dan dicatat untuk keterulangan.

{\footnotesize
\setlength{\tabcolsep}{4pt}
\renewcommand{\arraystretch}{1.15}
\begin{longtable}{|>{\centering\arraybackslash}p{0.04\textwidth}|L{0.24\textwidth}|L{0.30\textwidth}|L{0.32\textwidth}|}
\caption{Alat dan Bahan Penelitian}
\label{tab:tools}\\
\hline
\textbf{No.} & \textbf{Alat / Bahan} & \textbf{Spesifikasi / Keterangan} & \textbf{Fungsi dalam Penelitian} \\ \hline
\endfirsthead
\hline
\textbf{No.} & \textbf{Alat / Bahan} & \textbf{Spesifikasi / Keterangan} & \textbf{Fungsi dalam Penelitian} \\ \hline
\endhead
1 & Bahasa pemrograman & Python 3.11 atau lebih baru & \textit{Pipeline} standardisasi, analitik, dan analisis statistik \\ \hline
2 & Kerangka kerja \textit{backend} & FastAPI & API dan logika bisnis \\ \hline
3 & Kerangka kerja \textit{frontend} & Aplikasi \textit{web} responsif (React) & Formulir pendataan dan \textit{dashboard} \\ \hline
4 & Basis data & PostgreSQL dengan ekstensi pgvector & Penyimpanan data UMKM dan indeks vektor KBLI \\ \hline
5 & LLM & Satu model \textit{proprietary} melalui API dan satu model \textit{open-weight} lokal & Standardisasi deskripsi usaha \\ \hline
6 & Model \textit{embedding} & Model \textit{embedding} multibahasa yang mendukung Bahasa Indonesia & \textit{Retrieval} semantik kandidat KBLI \\ \hline
7 & Pustaka analitik & pandas, scikit-learn, scipy, statsmodels & \textit{Data profiling}, rekapitulasi, perhitungan metrik, dan uji statistik \\ \hline
8 & Alat anotasi & Label Studio & Penyusunan \textit{gold standard} KBLI \\ \hline
9 & Pemodelan proses & Notasi BPMN 2.0 (misalnya Bizagi Modeler atau Camunda Modeler) & Model proses \textit{as-is} dan \textit{to-be} \\ \hline
10 & Perangkat keras & Komputer dengan GPU untuk model \textit{open-weight}; \textit{server} aplikasi & Inferensi lokal dan penyebaran \textit{prototype} \\ \hline
11 & Kamus KBLI 2025 & Peraturan BPS Nomor 7 Tahun 2025 \cite{bps2025perban7} & Basis pengetahuan \textit{retrieval} \\ \hline
12 & Data UMKM dinas & Kiriman desa tahun 2026 atau kompilasi dinas dari kiriman yang sama & \textit{Gold standard}, \textit{data profiling}, dan kondisi sebelum \\ \hline
13 & Instrumen evaluasi & SUS dan \textit{raw} NASA-TLX & Evaluasi pengguna \\ \hline
14 & Pengelolaan versi & Git dan repositori daring & Pengelolaan kode, \textit{prompt}, dan dokumentasi \\ \hline
\end{longtable}
}

\section{Risiko dan Mitigasi}
\label{sec:risiko}

Risiko utama pelaksanaan penelitian beserta mitigasinya disajikan pada Tabel~\ref{tab:risiko}.

{\footnotesize
\setlength{\tabcolsep}{4pt}
\renewcommand{\arraystretch}{1.15}
\begin{longtable}{|>{\centering\arraybackslash}p{0.04\textwidth}|L{0.30\textwidth}|L{0.13\textwidth}|L{0.42\textwidth}|}
\caption{Risiko Penelitian dan Mitigasi}
\label{tab:risiko}\\
\hline
\textbf{No.} & \textbf{Risiko} & \textbf{Tingkat} & \textbf{Mitigasi} \\ \hline
\endfirsthead
\hline
\textbf{No.} & \textbf{Risiko} & \textbf{Tingkat} & \textbf{Mitigasi} \\ \hline
\endhead
1 & Keterlambatan akses data karena beban kerja staf dinas & Tinggi & Surat resmi segera setelah Seminar Proposal; permintaan data bertahap; titik keputusan G1 \\ \hline
2 & Partisipasi desa percontohan rendah atau izin pelaksanaan di desa belum terbit & Sedang & Pemilihan desa bersama dinas; pelatihan singkat di lokasi; surat pengantar dari dinas; jalur cadangan berupa studi terkontrol dengan staf dinas dan perwakilan desa pada data kiriman desa 2026; titik keputusan G3 \\ \hline
3 & Kinerja standardisasi tidak memenuhi kriteria & Sedang & Antrean verifikasi manusia; pembatasan tingkat granularitas; titik keputusan G2 \\ \hline
4 & Biaya API LLM melebihi anggaran atau kebijakan data membatasi penggunaan layanan eksternal & Sedang & Penggunaan model \textit{open-weight} lokal sebagai alternatif utama; pembatasan jumlah panggilan pada eksperimen \\ \hline
5 & Konektivitas internet di desa tidak stabil & Sedang & Antarmuka ringan dengan penyimpanan isian sementara; uji coba juga dilakukan di kantor desa atau kecamatan \\ \hline
6 & Jumlah operator desa hanya enam sampai delapan sehingga daya uji Wilcoxon terbatas & Tinggi & Satu sampai dua operator pada setiap desa di tiga sampai empat desa; desain \textit{within-subject}; RQ3 diperlakukan eksploratif dengan nilai $p$ eksak, ukuran efek, dan interval kepercayaan; apabila kurang dari enam, RQ3 dilaporkan secara deskriptif dengan ukuran efek tanpa kesimpulan inferensial; triangulasi dengan data kualitatif; pembatasan klaim generalisasi \\ \hline
7 & Perubahan versi atau ketersediaan model LLM selama penelitian & Rendah & Penguncian versi model dan dokumentasi tanggal akses; pelaporan hasil per versi \\ \hline
8 & Anotator kedua hanya mengodekan sekitar 30\% \textit{gold standard} sehingga sebagian label berlabel tunggal & Sedang & Himpunan bagian berstrata; kalibrasi pedoman pada 50 \textit{record}; konsultasi dengan staf dinas untuk kasus ambigu; pelaporan kappa pada himpunan bagian dan keterbatasannya \\ \hline
9 & Isian jenis usaha terlalu umum untuk dikodekan sampai tingkat kelompok lima digit & Sedang & Status \texttt{PERLU\_VERIFIKASI} dalam pedoman anotasi dan keluaran model; evaluasi pada beberapa tingkat granularitas; pelaporan proporsi \textit{record} yang tidak dapat dikodekan; titik keputusan G2 \\ \hline
10 & Tidak ada model \textit{open-weight} kandidat yang memenuhi ambang kelayakan, atau perangkat GPU dan biaya API tidak memadai untuk fase \textit{benchmarking} & Sedang & Ambang kelayakan ditetapkan bersama dinas sebelum pengukuran; kandidat dengan ukuran berbeda, termasuk model terkuantisasi; model \textit{proprietary} sebagai cadangan dengan masukan yang dibatasi pada atribut yang diizinkan (KNF-02); tindak lanjut mengikuti titik keputusan G2 \\ \hline
\end{longtable}
}

\section{Rencana Luaran dan Publikasi}
\label{sec:luaran}

Luaran penelitian meliputi: (1) \textit{prototype} sistem informasi pendataan dan pendukung keputusan pembinaan UMKM beserta dokumentasi; (2) \textit{gold standard} pengodean KBLI 2025 beranotasi (500 \textit{record}) dan pedoman anotasi; (3) \textit{pipeline} standardisasi yang terdokumentasi, termasuk hasil \textit{benchmarking} LLM, \textit{prompt}, konfigurasi model, dan skrip analisis; (4) laporan evaluasi eksperimen standardisasi, kualitas data, dan evaluasi pengguna; (5) laporan Tugas Akhir; serta (6) manuskrip artikel ilmiah untuk \textit{Journal of Computer Information Systems}.

Kesesuaian penelitian dengan cakupan JCIS disajikan pada Tabel~\ref{tab:jcis}. Manuskrip akan disusun sesuai ketentuan jurnal, yaitu judul tidak melebihi sepuluh kata, abstrak tidak terstruktur sekitar 150 kata, tiga sampai lima kata kunci, panjang naskah tidak melebihi 8.000 kata, serta pernyataan penggunaan kecerdasan buatan generatif dan pernyataan ketersediaan data. Judul kerja manuskrip adalah \textit{``Retrieval-Grounded LLMs for Standardizing MSME Data in Local Government''}. Kerangka manuskrip disusun pada tahap penulisan laporan.

{\footnotesize
\setlength{\tabcolsep}{4pt}
\renewcommand{\arraystretch}{1.15}
\begin{longtable}{|L{0.30\textwidth}|L{0.62\textwidth}|}
\caption{Kesesuaian Penelitian dengan Cakupan \textit{Journal of Computer Information Systems}}
\label{tab:jcis}\\
\hline
\textbf{Topik Cakupan JCIS} & \textbf{Komponen Penelitian} \\ \hline
\endfirsthead
\hline
\textbf{Topik Cakupan JCIS} & \textbf{Komponen Penelitian} \\ \hline
\endhead
\textit{Analytics, Business Intelligence, Decision Support Systems} & \textit{Dashboard} rekapitulasi dan pemeringkatan sederhana; kelanjutan agenda penelitian LLM untuk analisis teks \textit{business intelligence} \cite{lotfipoor2026llm} \\ \hline
\textit{Data Management, Data Mining, Database Design and Development} & Standardisasi data berbasis \textit{retrieval-grounded} LLM, pengukuran kualitas data pada enam dimensi \cite{iso25012,govuk2021dqdimensions}, dan perancangan basis data \\ \hline
\textit{Systems Analysis, Design, and/or Implementation} & Analisis proses \textit{as-is}/\textit{to-be}, rekayasa kebutuhan, dan pengembangan artefak dengan DSR \\ \hline
\textit{Human-Computer Interaction} & Evaluasi kinerja tugas, beban kerja, dan \textit{usability} oleh aparatur \\ \hline
\textit{Ethical and Legal Issues of IT} & Pelindungan data pribadi dan rancangan pemrosesan LLM yang membatasi pengiriman data pribadi \\ \hline
\end{longtable}
}

\section{Jadwal Penelitian}
\label{sec:jadwal}

Penelitian dilaksanakan dalam dua semester. Semester ganjil berlangsung pada 7 September--18 Desember 2026 dengan Seminar Proposal pada 9--13 November 2026 (minggu ke-10) dan Seminar Tugas Akhir 1 pada 14--18 Desember 2026 (minggu ke-15). Semester genap berlangsung pada 15 Februari--4 Juni 2027 dengan Seminar Hasil pada 19--23 April 2027 (minggu ke-10) dan Sidang Tugas Akhir pada 24--28 Mei 2027 (minggu ke-15). Jadwal disajikan per bulan dan per minggu pada Tabel~\ref{tab:jadwal1} dan Tabel~\ref{tab:jadwal2}, dengan kode penanggung jawab mengikuti Tabel~\ref{tab:peran}.

\begin{table}[H]
  \caption{Jadwal Pelaksanaan Semester Ganjil (7 September--18 Desember 2026)}
  \label{tab:jadwal1}
  \centering
  \adjustbox{max width=\textwidth}{%
  \footnotesize
  \setlength{\tabcolsep}{2.5pt}
    \begin{tabular}{|c|>{\RaggedRight\arraybackslash}p{5.2cm}|c|*{15}{>{\centering\arraybackslash}p{0.62cm}|}}
    \hline
    \multirow{3}{*}{\textbf{No.}} & \multirow{3}{*}{\textbf{Kegiatan}} & \multirow{3}{*}{\textbf{PJ}} & \multicolumn{15}{c|}{\textbf{Bulan}} \\ \cline{4-18}
     & & & \multicolumn{4}{c|}{\textbf{September}} & \multicolumn{4}{c|}{\textbf{Oktober}} & \multicolumn{5}{c|}{\textbf{November}} & \multicolumn{2}{c|}{\textbf{Desember}} \\ \cline{4-18}
     & & & 1 & 2 & 3 & 4 & 1 & 2 & 3 & 4 & 1 & 2 & 3 & 4 & 5 & 1 & 2 \\ \hline
    1 & Studi literatur dan analisis LKjIP serta Renstra & Semua & \ga & \ga & \ga & \ga & \ga & \ga & \ga & \ga &  &  &  &  &  &  &  \\ \hline
    2 & Wawancara pemangku kepentingan dan validasi masalah & Semua &  &  &  &  &  & \ga & \ga & \ga &  &  &  &  &  &  &  \\ \hline
    3 & Penyusunan proposal & Semua &  &  &  &  &  &  & \ga & \ga & \ga &  &  &  &  &  &  \\ \hline
    4 & \textbf{Seminar Proposal (9--13 Nov 2026)} & Semua &  &  &  &  &  &  &  &  &  & \ms &  &  &  &  &  \\ \hline
    5 & Pengurusan izin penelitian, akses data, dan izin uji coba desa & M1 &  &  &  &  &  &  &  & \ga & \ga & \ga &  &  &  &  &  \\ \hline
    6 & Pemodelan proses dan validasi kebutuhan & M1 &  &  &  &  &  &  &  &  & \ga & \ga & \ga & \ga &  &  &  \\ \hline
    7 & \textit{Data profiling} dan pengukuran kualitas data kiriman desa 2026 & M1 &  &  &  &  &  &  &  &  &  &  & \ga & \ga & \ga &  &  \\ \hline
    8 & Pedoman anotasi, kalibrasi, dan \textit{gold standard} pengembangan & M2 &  &  &  &  &  &  &  &  &  & \ga & \ga & \ga & \ga &  &  \\ \hline
    9 & Indeks KBLI, \textit{retriever}, dan eksperimen awal & M2 &  &  &  &  &  &  &  &  &  &  & \ga & \ga & \ga & \ga &  \\ \hline
    10 & \textit{Benchmarking} LLM pada data pengembangan (akurasi, biaya, latensi) dan pemilihan model & M2 &  &  &  &  &  &  &  &  &  &  &  & \ga & \ga & \ga &  \\ \hline
    11 & Perancangan arsitektur, basis data, dan antarmuka & M3 &  &  &  &  &  &  &  &  &  & \ga & \ga & \ga & \ga & \ga &  \\ \hline
    12 & Penetapan kriteria dan bobot pemeringkatan; rancangan instrumen evaluasi & M3 &  &  &  &  &  &  &  &  &  &  &  & \ga & \ga & \ga &  \\ \hline
    13 & Penyusunan laporan Tugas Akhir 1 & Semua &  &  &  &  &  &  &  &  &  &  &  &  & \ga & \ga &  \\ \hline
    14 & \textbf{Seminar Tugas Akhir 1 (14--18 Des 2026)} & Semua &  &  &  &  &  &  &  &  &  &  &  &  &  &  & \ms \\ \hline
    \end{tabular}%
  }
  \begin{flushleft}
    \small \textit{Keterangan:} angka di bawah nama bulan menunjukkan minggu dalam bulan (dihitung dari hari Senin); minggu ke-1 semester dimulai 7 September 2026. \colorbox{green!25}{\phantom{X}} Kegiatan penelitian \quad \colorbox{orange!50}{\phantom{X}} Seminar/\textit{milestone}
  \end{flushleft}
\end{table}

\begin{table}[H]
  \caption{Jadwal Pelaksanaan Semester Genap (15 Februari--4 Juni 2027)}
  \label{tab:jadwal2}
  \centering
  \adjustbox{max width=\textwidth}{%
  \footnotesize
  \setlength{\tabcolsep}{2.5pt}
    \begin{tabular}{|c|>{\RaggedRight\arraybackslash}p{5.2cm}|c|*{16}{>{\centering\arraybackslash}p{0.62cm}|}}
    \hline
    \multirow{3}{*}{\textbf{No.}} & \multirow{3}{*}{\textbf{Kegiatan}} & \multirow{3}{*}{\textbf{PJ}} & \multicolumn{16}{c|}{\textbf{Bulan}} \\ \cline{4-19}
     & & & \multicolumn{2}{c|}{\textbf{Februari}} & \multicolumn{5}{c|}{\textbf{Maret}} & \multicolumn{4}{c|}{\textbf{April}} & \multicolumn{5}{c|}{\textbf{Mei}} \\ \cline{4-19}
     & & & 3 & 4 & 1 & 2 & 3 & 4 & 5 & 1 & 2 & 3 & 4 & 1 & 2 & 3 & 4 & 5 \\ \hline
    1 & Penyelesaian \textit{gold standard} uji (400 \textit{record}) & M2 & \ga & \ga & \ga &  &  &  &  &  &  &  &  &  &  &  &  &  \\ \hline
    2 & Eksperimen standardisasi B0--P2b dan analisis & M2 &  &  &  & \ga & \ga & \ga & \ga & \ga &  &  &  &  &  &  &  &  \\ \hline
    3 & Implementasi modul pendataan dan validasi & M1 & \ga & \ga & \ga & \ga & \ga &  &  &  &  &  &  &  &  &  &  &  \\ \hline
    4 & Implementasi \textit{dashboard} rekapitulasi dan pemeringkatan sederhana & M3 &  & \ga & \ga & \ga & \ga & \ga &  &  &  &  &  &  &  &  &  &  \\ \hline
    5 & Integrasi layanan standardisasi dan antrean verifikasi & M2, M3 &  &  &  & \ga & \ga & \ga &  &  &  &  &  &  &  &  &  &  \\ \hline
    6 & \textit{Functional testing} dan pembekuan fitur & Semua &  &  &  &  & \ga & \ga &  &  &  &  &  &  &  &  &  &  \\ \hline
    7 & Uji coba lapangan tiga minggu pada 3--4 desa percontohan & M1 &  &  &  &  &  &  & \ga & \ga & \ga &  &  &  &  &  &  &  \\ \hline
    8 & Evaluasi \textit{within-subject} (tugas terkontrol, NASA-TLX, SUS, wawancara) & M3 &  &  &  &  &  &  & \ga & \ga & \ga &  &  &  &  &  &  &  \\ \hline
    9 & Pengukuran kualitas data sesudah dan analisis & M1 &  &  &  &  &  &  &  &  & \ga & \ga & \ga &  &  &  &  &  \\ \hline
    10 & \textbf{Seminar Hasil (19--23 Apr 2027)} & Semua &  &  &  &  &  &  &  &  &  & \ms &  &  &  &  &  &  \\ \hline
    11 & Analisis lanjutan dan sintesis hasil & Semua &  &  &  &  &  &  &  &  &  &  & \ga & \ga &  &  &  &  \\ \hline
    12 & Penyusunan laporan Tugas Akhir & Semua &  &  &  &  &  &  &  &  & \ga & \ga & \ga & \ga & \ga & \ga &  &  \\ \hline
    13 & Penyusunan manuskrip JCIS & Semua &  &  &  &  &  &  &  &  &  &  & \ga & \ga & \ga & \ga &  & \ga \\ \hline
    14 & \textbf{Sidang Tugas Akhir (24--28 Mei 2027)} & Semua &  &  &  &  &  &  &  &  &  &  &  &  &  &  & \ms &  \\ \hline
    15 & Revisi laporan dan pengajuan manuskrip & Semua &  &  &  &  &  &  &  &  &  &  &  &  &  &  &  & \ga \\ \hline
    \end{tabular}%
  }
  \begin{flushleft}
    \small \textit{Keterangan:} angka di bawah nama bulan menunjukkan minggu dalam bulan (dihitung dari hari Senin); minggu ke-1 semester dimulai 15 Februari 2027 dan minggu ke-16 berakhir 4 Juni 2027. \colorbox{green!25}{\phantom{X}} Kegiatan penelitian \quad \colorbox{orange!50}{\phantom{X}} Seminar/\textit{milestone}
  \end{flushleft}
\end{table}

Pelaksanaan penelitian dikendalikan melalui empat titik keputusan (\textit{decision gate}) sebagaimana Tabel~\ref{tab:gate}. Setiap titik keputusan memiliki kriteria yang terukur dan rencana alternatif apabila kriteria tidak terpenuhi.

{\footnotesize
\setlength{\tabcolsep}{4pt}
\renewcommand{\arraystretch}{1.15}
\begin{longtable}{|>{\centering\arraybackslash}p{0.06\textwidth}|L{0.16\textwidth}|L{0.33\textwidth}|L{0.34\textwidth}|}
\caption{Titik Keputusan Penelitian}
\label{tab:gate}\\
\hline
\textbf{Kode} & \textbf{Waktu} & \textbf{Kriteria} & \textbf{Rencana Alternatif} \\ \hline
\endfirsthead
\hline
\textbf{Kode} & \textbf{Waktu} & \textbf{Kriteria} & \textbf{Rencana Alternatif} \\ \hline
\endhead
G1 & Minggu ke-12 semester ganjil (23--27 Nov 2026) & Data dinas diterima dan izin penelitian terbit & Penyusunan \textit{gold standard} dimulai dari kiriman desa 2026 yang telah tersedia; pengajuan ulang melalui pimpinan program studi \\ \hline
G2 & Seminar TA-1 (14--18 Des 2026) & Fase \textit{benchmarking} LLM selesai, dan model utama (\textit{open-weight}) serta model replikasi (\textit{proprietary}) tercatat pada \textit{log}; Acc@1 golongan pokok kondisi P1 pada data pengembangan (100 \textit{record}) lebih tinggi daripada kondisi B0 dan B1, dibandingkan secara deskriptif tanpa uji signifikansi karena sampel kecil & Apabila tidak ada kandidat yang memenuhi ambang kelayakan, ambang ditinjau bersama dinas atau kandidat model lain diuji; perbaikan \textit{retriever}, glosarium, dan contoh \textit{few-shot} pada data pengembangan; apabila tetap tidak terpenuhi, seluruh keluaran diteruskan ke antrean verifikasi manusia dan hasil dilaporkan apa adanya \\ \hline
G3 & Minggu ke-6 semester genap (22--26 Mar 2027) & \textit{Prototype} lulus \textit{functional testing}; sekurang-kurangnya tiga desa percontohan bersedia berpartisipasi dan izin pelaksanaan terbit & Studi terkontrol dengan staf dinas dan perwakilan desa pada data kiriman desa 2026 \\ \hline
G4 & Minggu ke-9 semester genap (12--16 Apr 2027) & Data evaluasi RQ2 dan RQ3 lengkap & Seminar Hasil menyajikan hasil RQ1 secara lengkap dan hasil awal RQ2--RQ3; pengumpulan data dilanjutkan hingga minggu ke-11 \\ \hline
\end{longtable}
}